% Modernised reading — fonds Alexandre Grothendieck, shelfmark 83
% (pages 1-221, the whole folder).
%
% This is an INTERPRETATION, not a transcription. It re-reads the pages in
% today's notation and vocabulary, states the hypotheses the manuscript leaves
% implicit, and drops the critical apparatus so that the mathematics reads
% straight through. Everything the brackets and underlines carried moves into a
% footnote. It is held to being mathematically correct as it stands; where the
% manuscript is loose or mistaken, the true statement is what appears here and
% the footnote says what the page has. In any dispute of substance the
% transcription governs: whoever cites Grothendieck must cite batch-01.fr.tex
% to batch-12.fr.tex (pp. 1-20, 21-40, ..., 201-220, 221).
%
% Pass: Opus 5.5 (claude-opus-5-5), 2026-10-17 — first pass, unchecked against the pages by a human.
% The folder was read whole, its twelve batches together; this is one document
% for the shelfmark. It was made on Opus 5.5 and has not been measured against
% the reference reading of folder 115 (made on Opus 5). The literature named in
% the footnotes is cited from memory and was not checked against the sources.
% Re-done symbolically or numerically for this reading: the counts and traces
% over F_4 (pp. 2-15); the braid computations (pp. 18-20); the SL_2 identities
% (pp. 21-36); the multiplication table, trace matrix, discriminants and the
% invariant rho_0 of M(b,c,b',c',t) (pp. 83-99, 116-120); t_{phi,phi'} as a
% root of Q (pp. 138-139); Sigma, Pi, Delta for three matrices (pp. 183-188);
% the delta relation on A*/k* (p. 77); the two determinants of pp. 210, 218.
%
% Not transcribed, hence not read: p. 1 (blank); pp. 4, 6, 8, 10, 12, 14, 16,
% 93, 195, 197, 201-219 odd (typescripts of his non-mathematical writing);
% covers pp. 17, 37, 51, 74 (their words are used as titles); blanks pp. 46,
% 100, 148, 181, 205; pp. 76, 78, 79, 128 (no note in the transcription).
%
% Spine. Combinatorial polyhedra (the icosahedron, the congruence maps of
% folder 79) realised inside the algebra of 2x2 matrices, then that algebra
% (a rank-4 Azumaya algebra over a ring) studied through the invariants of
% one, two, three, four of its elements:
%   pp. 2-15    the icosahedron as the 12 u in M_2(F_4) with Tr u = 1,
%               det u = alpha; Tr(uv) separates edges and co-edges.
%   pp. 18-20   the braid relation, B_3 and SL_2(Z) (Nov. 1982, 19.6.83).
%   pp. 21-28   SL_2 scratch, and stray leaves of the axiomatics of folder 79.
%   pp. 29-36   « Structure du groupe SL(2, Lambda) », author's pp. 1-4.
%   pp. 38-50   « Rapport anharmonique » and n-alities (April 1986).
%   pp. 52-60   « Bi-icosaèdres et cartes de congruences » (April 1986).
%   pp. 61-73   the nilpotent cone of a rank-4 Azumaya algebra, its classes
%               C*(e), the oriented triangle structure.
%   pp. 75-77   attempts: the invariant delta on A*/k*.
%   pp. 80-99   Azumaya rang 4 (I): the algebra M(b,c,b',c';t), rho, Q.
%   pp. 101-147 Azumaya rang 4 (II): rank-3 submodules, two quadratic
%               subalgebras, rho_0, the intrinsic theory, X wedge X'.
%   pp. 149-180 Azumaya rang 4 (III): affine spaces with involution, the
%               product *, lambda-split extensions.
%   pp. 182-200 (I bis), May-July 1986: Tr xyz, Sigma, Pi, Delta; four
%               variables and the cube.
%   pp. 202-220 linear relations on the edges of the cube.
%   p. 221      a typescript page (L 19) of his reflection on the milieu.
%
% Notation fixed for the whole folder. k a commutative ring; A an Azumaya
% algebra of rank 4 over k (the pages write A or M); Tr, det its reduced trace
% and norm; u' = Tr(u).1 - u its canonical involution; phi(u,v) = Tr(uv) and
% psi(u,v) = Tr u Tr v - Tr(uv), the polar form of det; A^nat = ker Tr.
% For two elements u, v: b = Tr u, c = det u, b' = Tr v, c' = det v,
% t = Tr(uv); delta = b^2 - 4c, delta' = b'^2 - 4c';
% Q = t^2 - bb't + b^2c' + b'^2c - 4cc'; rho_0 = rho - 2 = delta delta'/Q.
% For three: tau = Tr xyz, tau' = Tr zyx, Sigma, Pi, Delta as on pp. 183-188.
% The letter rho names five different things on the pages (braid elements
% p. 18; rho(beta) in SL_2 p. 34; cross-ratios rho_ijkl p. 42; the invariant
% rho = lambda + lambda^{-1} p. 86; an element of order 3 p. 69). Only the
% fourth is used outside its section; the fifth is renamed r here. The
% commutator written x*y on p. 77 is written [x,y] here, the star being kept
% for the product of quadratic algebras and affine spaces (pp. 142-180). As on
% the pages and in folder 79, E* (for E free of rank 2) is the set of vectors
% that are part of a basis, L* the generators of an invertible L.
%
% Substantive departures from the pages, with page numbers:
%  - pp. 2, 5, 11: det(u+v) over F_4 lacks the term Tr u Tr v; the correct
%    det(u+v) = 1 + Tr(uv) agrees with the table of p. 15, which is right.
%  - p. 18: c)-e) hold in SL_2(Z), not in B_3 (the margin says so); in B_3,
%    rho_0^3 = sigma^{-2} and rho_0' = sigma^2 rho_0. p. 20 writes rho_0^2
%    for rho_0^{-2}.
%  - p. 27: the boxed product formula has lambda'^{-2} where lambda'^2 is right
%    (twice).
%  - p. 31, (15): e_0 for e_1; (17): beta in L_0, not Lambda.
%  - pp. 61-62: « rho_f d'ordre 6 » is order 3 in SL_2 (6 up to sign).
%  - p. 69: Tr r = -1 for the order-3 r, not 1; then r^{-1} = -1 - r and
%    tau = ±(1 + 2r); the page's formulas are those of -r.
%  - p. 77: delta(xy) + delta(xy^{-1}) = delta x delta y + 2(delta x +
%    delta y) holds only when Tr(xyx^{-1}y^{-1}) = 2 (the page doubts it); the
%    general correction term is given; the polynomial P of the commutator
%    does not exist as a polynomial in delta(x), delta(y), delta(xy).
%  - p. 80, (2): Hamilton-Cayley reads « = 1 » for « = 0 ».
%  - p. 86: « P - 2Q ≠ 2 » read as « P - 2Q = delta delta' invertible ».
%  - p. 92: Tr((uv)^2) = t^2 - 2cc', not t^2 - cc'.
%  - p. 96: Tr(ux) = 0 for nilpotent u means x in the Borel of u, not « [x,u]
%    = 0 »; p. 101-102: the two discriminants agree on B ∋ 1 of odd rank,
%    and differ by (-1)^{rk B - 1} in general (the page says « égales » on V,
%    where they are opposite). Hence p. 118's box needs a sign:
%    rho_0 = - delta_H delta_H' / delta_B (p. 120 has -delta_B = Q).
%  - p. 143: det(u*v) = b^2c' + b'^2c - 4cc'.
%  - p. 188: Pi lacks the term -4cc'c''; the page's boxed Pi is false when
%    cc'c'' ≠ 0. p. 184: « 2(tau - tau') = eps_0 Delta » is tau - tau' =
%    eps_0 Delta, as p. 188 has.
%  - p. 218: the 6x6 determinant is 512 = 2^9, not 27^2 - 1 = 728; the
%    conclusion (non-zero, prime to 3) stands.
%
% Announced and never established in the folder: the end of the braid
% computation (p. 20); (70) on p. 24; the use of N' = T ⊔ N^- (p. 36); the
% verification (7) and the proposition of pp. 43-44 on quadralities; the
% simpler description of the category of quadralities (p. 45); the explicit
% lambda-orientation of the classes C*(e) (p. 67); which of r, r^{-1} the
% relation tau = ±(1+2r) selects (p. 70); the polynomial Q of p. 77; the
% sufficiency half of the theorem of p. 98 (p. 99 stops; supplied here when 2
% is invertible); (iii) => (i) of p. 108 over a non-reduced base (doubted
% p. 106); Zariski-local realisability of rho_0 (p. 115); multiple contracted
% products and their quadratic conditions (p. 147); Q(A_1 * A_2) ≃
% Q(A_1) * Q(A_2) (pp. 168-170); the unit for *_lambda (pp. 179-180); the
% freeness of B over A in four variables and the relation among the 14
% invariants (pp. 192-193); the 32 relations of the cube (p. 206).
\documentclass[11pt,a4paper]{article}
\input{../preamble/grothendieck}

\folder{83}
\pages{1}{221}
\dating{1982-1986}
\foldertitle{[Icosaèdres, bi-icosaèdres et algèbres d'Azumaya de rang 4] : notes manuscrites (1982-1983, 1986, s.d.) — lecture modernisée du dossier entier}

\begin{document}

\section*{Résumé}

\begin{resume}
Ces feuillets prennent l'une des plus vieilles figures de la géométrie,
l'icosaèdre — vingt triangles, douze sommets, cinq triangles autour de chaque
sommet — et la cherchent là où on ne l'attend pas : parmi les tableaux de
nombres $2 \times 2$. Sur un corps à quatre éléments, il y a exactement douze
matrices de trace $1$ et d'un déterminant convenablement choisi ; on décrète que deux d'entre
elles sont reliées par une arête quand la trace de leur produit prend une
certaine valeur, et l'on obtient un icosaèdre. Le même dessin sort des
matrices nilpotentes à coefficients modulo $5$, et plus généralement, modulo
$n$, des « cartes de congruence » qu'un dossier voisin (le dossier 79) avait
construites à partir des vecteurs du plan. Le dossier est l'histoire de ce
déplacement : on part d'un polyèdre, on aboutit à l'étude de l'algèbre des
matrices $2 \times 2$ elle-même.

Le problème est de comprendre cette algèbre sans coordonnées, sur une base
quelconque — sur les entiers, sur un anneau où $2$ n'est pas inversible —, là
où elle s'appelle algèbre d'Azumaya de rang $4$ (une algèbre de quaternions,
dit-on aussi). Deux nombres décrivent un élément : sa trace et son
déterminant. Pour deux éléments, il en faut un de plus, la trace du produit, et
le dossier montre que ces cinq nombres suffisent à reconstruire l'algèbre
qu'ils engendrent, et quand elle est « toute » l'algèbre. Le critère fait
apparaître une quantité qui joue le rôle d'un rapport anharmonique — le
nombre qui, en géométrie projective, mesure la position mutuelle de quatre
points sur une droite. Ici les quatre points sont les deux paires de points
fixes de deux « rotations » ; leur rapport anharmonique est la seule chose
qui ne change pas quand on fait tourner le tout. Pour trois éléments, la trace
du produit $xyz$ n'est plus déterminée par les traces et déterminants : elle
est l'une des deux racines d'une équation du second degré, et l'autre racine
est la trace de $zyx$. L'ordre dans lequel on multiplie, c'est-à-dire
l'orientation, est exactement ce qui manque.

Une idée court sous tout cela, et elle prend dans le dossier une forme
générale. Une algèbre « quadratique » (comme les nombres complexes au-dessus
des réels) porte une symétrie, la conjugaison. Pour en combiner deux, on ne
prend pas leur produit tensoriel, trop gros, mais la partie de ce produit que
les deux conjugaisons, faites ensemble, laissent fixe. Quand $2$ est
inversible, ce n'est qu'une manière savante de multiplier deux
discriminants ; en caractéristique $2$, c'est une opération délicate, et une
quarantaine de pages construisent, pour la faire marcher sans hypothèse, une
algèbre linéaire d'« espaces affines à involution ». Le dossier voisin 84
reprend cette construction plus à fond.

Le dossier ne conclut pas. Il s'arrête sur des relations entre traces de
produits de quatre matrices, disposées sur les sommets, arêtes et faces d'un
cube, avec en marge « Trop lourd ». Les dates (novembre 1982, juin 1983, puis
avril à juillet 1986) le placent à côté de la rédaction de \emph{Récoltes et
Semailles}, dont des pages dactylographiées servent de brouillon au verso. Le
dernier feuillet est l'une d'elles : il y décrit un « aplatissement » de la
pensée mathématique (le mot est, dit-il, de Siegel), le mépris de son
« versant féminin » — l'inspiration, le rêve, la vision, et même l'idée tant
qu'elle n'est pas un énoncé publié. Les pages de calcul qui précèdent en sont
comme l'envers : on y voit une figure venir avant les énoncés, des « on
gagne ! », un « C'est idiot ! » en marge d'une conjecture, des vérifications
remises à plus tard et des théorèmes notés « sauf erreur ».

Les noms à chercher ensuite : groupe $\mathrm{PSL}_2(\mathbf{F}_4) \simeq
\mathfrak{A}_5$, grand icosaèdre, courbe modulaire $X(5)$, groupe de tresses
$B_3$ et $\mathrm{SL}_2(\mathbf{Z})$, relations de Steinberg, rapport
anharmonique, carrés latins et réseaux, algèbres de quaternions et cône
nilpotent, produit des algèbres quadratiques, identités de traces de Fricke,
invariants de matrices $2 \times 2$.

\keywords{icosahedron, PSL2(F4), A5, great icosahedron, braid group B3, SL2(Z), Steinberg relations, Gauss decomposition, transporter, cross-ratio, Latin square, net, transversal design, bitorsor, congruence map, modular curve X(5), PGL2(F5) = S5, nilpotent cone, Azumaya algebra, quaternion algebra, Cayley-Hamilton, reduced norm, quaternion algebra generated by two elements, Borel subalgebra, relative position of two tori, quadratic algebra, quadratic cover, contracted product, product of quadratic algebras, discriminant, affine space with involution, Fricke trace identities, invariants of 2x2 matrices, trace ring}
\end{resume}

\pagerange{2}{221}
\section*{Le fil du dossier, et les conventions}

Le dossier n'a pas de pagination continue de sa main. Il est fait de suites
distinctes, dont plusieurs sont ouvertes par une couverture portant un titre
et une date : « Rapport anharmonique, avril 1986 » (page 37), « Bi-icosaèdres
et cartes de congruences, avril 1986 » (page 51), « Azumaya rang 4, avril
1986 » (page 74) ; puis des titres de sa main en tête de page : « Azumaya
rang 4 (I) » (page 80), « (II) » (page 101), « (III) » (page 149), « (I bis) »
(page 182, mai--juillet 1986). Les pages 2 à 20 sont plus anciennes
(novembre 1982 et 19 juin 1983 aux pages 18 et 19). Des pages dactylographiées
de ses textes non mathématiques sont mêlées aux feuillets ; elles ne sont pas
transcrites, sauf la dernière (page 221), qui porte sa main.

\subsection*{Les stations}

\begin{itemize}
  \item \emph{Pages 2 à 15 : l'icosaèdre sur $\mathbf{F}_4$.} Les douze
    matrices de trace $1$ et de déterminant $\alpha$ ; la trace du produit
    distingue le même sommet, l'antipode, l'arête et la « coarête » ; les
    étoiles pentagonales.
  \item \emph{Pages 18 à 20 : la relation de tresses.} Le groupe
    $\langle \varepsilon_0, \varepsilon_1 \mid \varepsilon_0\varepsilon_1\varepsilon_0
    = \varepsilon_1\varepsilon_0\varepsilon_1\rangle$, ses éléments $\sigma$,
    $\rho_i$, $\rho_i'$, et son quotient $\mathrm{SL}_2(\mathbf{Z})$.
  \item \emph{Pages 21 à 28 : brouillons.} Calculs dans $\mathrm{SL}_2$ sur un
    anneau, mêlés de feuillets d'une autre suite (pages 22, 24, 26, 28) qui
    parlent la langue du § XI du dossier 79.
  \item \emph{Pages 29 à 36 : « Structure du groupe $\mathrm{SL}(2,
    \Lambda)$ ».} Tore déployé, unipotents, Borels, l'application $\sigma_0$,
    transporteurs, et la caractérisation de l'accouplement $L_0 \otimes L_1
    \simeq \Lambda$ par la relation de tresses.
  \item \emph{Pages 38 à 50 : rapport anharmonique et $n$-alités.} Trois
    puis quatre flèches de même source dont deux quelconques forment un
    produit ; les rapports anharmoniques qui en sortent ; la version
    ensembliste (structures d'incidence).
  \item \emph{Pages 52 à 60 : bi-icosaèdres et cartes de congruence.} Les deux
    icosaèdres complémentaires sur douze sommets, leurs orientations, leurs
    groupes ; la réalisation des cartes de congruence dans le cône nilpotent de
    $M_2(\mathbf{Z}/n)$.
  \item \emph{Pages 61 à 73 : le cône nilpotent d'une algèbre d'Azumaya.} Ses
    classes $C^{*}(e)$, la structure de graphe $\operatorname{Tr}(uv) = -1$, les
    triangles et leur orientation canonique.
  \item \emph{Pages 75 à 77 : tentatives.} Décrire $M_2(k)$ par le groupe
    $A^{*}/k^{*}$ et la fonction $\delta(x) = (b^2 - 4c)/c$.
  \item \emph{Pages 80 à 99 : Azumaya rang 4 (I).} Hamilton--Cayley ;
    l'algèbre $M(b, c; b', c'; t)$ ; l'invariant anharmonique ; le théorème de
    la page 98.
  \item \emph{Pages 101 à 147 : Azumaya rang 4 (II).} Sous-modules de rang
    $3$ contenant $1$ ; deux sous-algèbres quadratiques en position générale et
    l'invariant $\rho_0$ ; la réalisation ; la théorie intrinsèque
    (« twists ») ; le produit contracté $X \wedge X'$ des revêtements
    quadratiques.
  \item \emph{Pages 149 à 180 : Azumaya rang 4 (III).} Produits d'espaces
    affines à involution ; le torseur $Q(A)$ ; extensions $\lambda$-scindées
    et leur produit.
  \item \emph{Pages 182 à 200 : (I bis).} La trace de $xyz$ ; l'anneau des
    invariants de trois matrices ; quatre matrices et le cube.
  \item \emph{Pages 202 à 220 : le cube.} Relations linéaires entre quantités
    portées par les arêtes ; paramétrage ; deux déterminants.
  \item \emph{Page 221 : un feuillet dactylographié} (« L 19 ») de sa
    réflexion sur le milieu mathématique.
\end{itemize}

Ce ne sont pas les chapitres d'un exposé. Les suites (I), (II), (III) et
(I bis) d'avril à juillet 1986 forment le seul fil continu ; les stations
antérieures en sont les sources. La « Structure du groupe $\mathrm{SL}(2,
\Lambda)$ » (pages 29 à 36) est une mise au propre des brouillons des pages
21 à 28, et ces deux stations restent distinctes.

\subsection*{Conventions, valables pour tout le dossier}

$k$ est un anneau commutatif (les pages disent aussi $\Lambda$ quand il est
le coefficient d'un module de rang $2$). $A$ est une algèbre d'Azumaya de rang
$4$ sur $k$, c'est-à-dire localement isomorphe à $M_2(k)$ ; les pages
l'appellent tour à tour $A$, $M$ ou $\mathcal{A}$. On note
$\operatorname{Tr}$ et $\det$ sa trace et sa norme réduites, $u' =
\operatorname{Tr}(u)\,1 - u$ son involution canonique (de sorte que $uu' =
\det(u)\,1$), et l'on a la relation de Hamilton--Cayley $u^2 -
\operatorname{Tr}(u)\,u + \det(u) = 0$. Deux formes bilinéaires symétriques
reviennent partout :
\[
  \varphi(u, v) = \operatorname{Tr}(uv), \qquad
  \psi(u, v) = \operatorname{Tr}u\,\operatorname{Tr}v - \operatorname{Tr}(uv),
\]
la seconde étant la forme polaire de $\det$ : $\det(u+v) = \det u + \det v +
\psi(u, v)$ et $\psi(u, u) = 2\det u$. On note $A^{\natural} =
\{u \mid \operatorname{Tr}u = 0\}$, orthogonal de $1$ pour l'une et l'autre
forme.

Pour deux éléments $u$, $v$, on pose, comme la page 83,
\[
  b = \operatorname{Tr}u,\quad c = \det u,\quad b' = \operatorname{Tr}v,\quad
  c' = \det v,\quad t = \operatorname{Tr}(uv),
\]
\[
  \delta = b^2 - 4c,\qquad \delta' = b'^2 - 4c',\qquad
  Q = t^2 - bb't + b^2c' + b'^2c - 4cc',
\]
et $\rho_0 = \rho - 2 = \delta\delta'/Q$ (l'invariant anharmonique
« rectifié »). Pour trois éléments $x, y, z$, la page 183 renomme : $b, b',
b''$ les traces, $c, c', c''$ les déterminants, $t = \operatorname{Tr}yz$, $t'
= \operatorname{Tr}zx$, $t'' = \operatorname{Tr}xy$ ; ce renommage ne vaut que
dans la station (I bis).

Pour un module $E$ libre de rang $2$, $E^{*}$ désigne, comme au dossier 79,
l'ensemble des vecteurs qui font partie d'une base, et $L^{*}$ l'ensemble des
générateurs d'un module inversible $L$ — non un dual\footnote{Le dual est noté
$\check{E}$ ou $E^{\vee}$. Les pages écrivent aussi $\mathrm{Gl}$, $\mathrm{Sl}$,
$\mathrm{Se}$ ; on écrit $\mathrm{GL}$, $\mathrm{SL}$.}.

\emph{Noms partagés.} La lettre $\rho$ désigne sur les pages cinq objets :
des éléments du groupe de tresses (page 18), l'élément $\rho(\beta)$ de
$\mathrm{SL}_2$ (page 34), les rapports anharmoniques $\rho_{ijkl}$ d'une
quadralité (page 42), l'invariant $\rho = \lambda + \lambda^{-1}$ de deux
sous-algèbres quadratiques (page 86 et suivantes), et un élément d'ordre $3$
(page 69). Seul le quatrième sert hors de sa station ; le cinquième est noté
$r$ ici. De même $\delta$ est le discriminant $b^2 - 4c$ à partir de la page
84, mais la fonction $(b^2 - 4c)/c$ sur $A^{*}/k^{*}$ à la page 77, et des
paramètres sur le cube aux pages 212 à 214. L'étoile $*$ est réservée au
produit des algèbres quadratiques et des espaces affines à involution (pages
142 à 180) ; le commutateur, qu'il note $x * y$ à la page 77, est écrit
$[x, y]$.

\subsection*{Ce que seule une lecture d'ensemble peut dire}

\emph{Une seule figure, trois réalisations.} L'icosaèdre des pages 2 à 15
(traces dans $M_2(\mathbf{F}_4)$), celui des pages 52 à 60 (vecteurs de
$(\mathbf{Z}/5)^2$ au signe près, comme au dossier 79) et celui des pages 61 à
73 (nilpotents de $M_2(\mathbf{F}_5)$) sont le même objet vu trois fois, et
les pages le savent sans le dire en une phrase : la relation d'arête est
chaque fois une valeur prescrite de $\operatorname{Tr}(uv)$ — $\alpha$ sur
$\mathbf{F}_4$ (page 11), $-1$ pour les nilpotents (page 73). Les
« coarêtes » de la page 11, où $\operatorname{Tr}(uv) = \alpha'$, sont les arêtes
de l'icosaèdre complémentaire des pages 54 et 55 : le « bi-icosaèdre » est déjà
là en 1982\footnote{Rapprochement de nous. Sur $\mathbf{F}_4$, le Frobenius
$x \mapsto x^2$ échange $\alpha$ et $\alpha'$, donc les deux icosaèdres ; le
groupe qu'on obtient en lui adjoignant $\mathrm{PSL}_2(\mathbf{F}_4)$ est
$\mathrm{P\Gamma L}_2(\mathbf{F}_4) \simeq \mathfrak{S}_5 \simeq
\mathrm{PGL}_2(\mathbf{F}_5)$, le groupe de la page 58.}.

\emph{La relation de tresses relie trois stations.} La relation
$\varepsilon_0\varepsilon_1\varepsilon_0 = \varepsilon_1\varepsilon_0\varepsilon_1$
des pages 18 à 20, la formule (32) de la page 34, $e_1(\gamma)e_0(\beta)e_1(\gamma)
= \sigma_0(\beta)$ pour $\beta\gamma = -1$, et la caractérisation (38) de la
page 35 sont un seul fait : dans $\mathrm{SL}_2(\Lambda)$, deux unipotents
opposés satisfont la relation de tresses exactement quand leurs paramètres ont
pour produit $-1$, et c'est ce qui fait de $\mathrm{SL}_2(\mathbf{Z})$ un
quotient du groupe de tresses $B_3$, avec $\sigma \mapsto \sigma_0(1)$. La
« carte universelle » des pages 62 et 63, engendrée par un élément d'ordre $4$
et un élément d'ordre $3$, en est la version cartographique.

\emph{Deux rapports anharmoniques.} Celui des pages 38 à 45 est un
automorphisme $\rho_{ijkl}$ d'un objet $V_i$ (un scalaire quand les
$\operatorname{End}(V_i)$ sont commutatifs) ; celui des pages 86 et 111 est le
birapport des deux paires de points fixes de deux sous-algèbres quadratiques
sur la conique. Les deux sont datés d'avril 1986 ; le second peut se lire comme un cas
particulier du premier (quatre droites isotropes du plan de $A^{\natural}$), mais
les pages ne font pas le lien, qui est de nous.

\emph{Le dossier 84.} Les pages 126 à 180 développent la même théorie que la
deuxième partie du dossier 84 (« $\chi$-scindages et produits contractés »,
pages 42 à 121 de ce dossier, en particulier la « digression affine » de ses
pages 65 à 79, et les $\lambda$-rétractions de ses pages 133 à 139). Les deux
dossiers sont datés de la même période ; on ne sait lequel précède.

\subsection*{Ce que le dossier annonce sans l'établir}

La fin du calcul de la page 20 ; la relation (70) de la page 24 ; l'usage de
$N' = T \sqcup N^{-}$ (page 36, où la rédaction s'arrête) ; la vérification
de la relation (7) et la proposition des pages 43 et 44 sur les quadralités ;
la description simple de leur catégorie (page 45) ; la
« $\lambda$-orientation » des classes $C^{*}(e)$ (page 67) ; le choix entre
$r$ et $r^{-1}$ (page 70) ; le polynôme $Q$ de la page 77 ; la suffisance du
théorème de la page 98 ; l'implication (iii) $\Rightarrow$ (i) de la page 108
sur une base non réduite ; la réalisabilité Zariski-locale de $\rho_0$ (page
115) ; les produits contractés multiples (page 147) ; l'isomorphisme
$Q(A_1 * A_2) \simeq Q(A_1) * Q(A_2)$ (pages 168 à 170) ; l'unité de
$*_\lambda$ (page 180) ; la structure de l'anneau des invariants de quatre
matrices et la relation entre ses quatorze générateurs (pages 192 et 193) ; les
trente-deux relations du cube (page 206).

\pagerange{2}{15}
\section*{I. L'icosaèdre sur $\mathbf{F}_4$ (pages 2 à 15)}

\pagerange{2}{7}
\subsection*{Les douze matrices (pages 2 à 7)}

Soit $k = \mathbf{F}_4$, $\alpha$ une racine de $\alpha^2 + \alpha + 1 = 0$ et
$\alpha' = \alpha^2 = \alpha + 1$ l'autre ; on calcule en caractéristique $2$.
Soit $V = k^2$, $M = \operatorname{End}(V)$, et
\[
  S = \{u \in M \mid \operatorname{Tr}u = 1,\ \det u = \alpha\}.
\]
Le polynôme caractéristique $T^2 + T + \alpha$ est irréductible sur $k$, et la
page 5 compte les $u$ de trace $1$ selon leur déterminant :
\[
  \det u = 0 : 20 \text{ (projecteurs de rang } 1),\quad
  \det u = 1 : 20,\quad
  \det u = \alpha : 12,\quad
  \det u = \alpha' : 12,
\]
soit $64 = 4^3$ en tout\footnote{Ce sont les nombres classiques $q^2 + q$
(polynôme scindé à racines distinctes) et $q^2 - q$ (polynôme irréductible)
pour $q = 4$. La page note d'abord, au lieu de « valeurs propres $\alpha$,
$\alpha'$ » pour $\det u = 1$, « $u = \mathrm{id} + $ nilpotent », qu'elle biffe.}.
L'ensemble $S$ a donc douze éléments, et le groupe $M^{*}/k^{*} =
\mathrm{PGL}_2(\mathbf{F}_4) = \mathrm{PSL}_2(\mathbf{F}_4) \simeq
\mathfrak{A}_5$ y opère par conjugaison ; la marge de la page 7 note qu'il y
est transitif\footnote{« Cor. 1. $M^{*}/k^{*}$ transitif — Cor. 2. $G$ sur $S$
… transitif », en marge, avec un mot illisible. La transitivité vient de ce
que deux éléments de même polynôme caractéristique irréductible sont
conjugués.}.

Pour $u \in S$ on a $u^2 = u + \alpha$, $u^4 = u + 1 = u'$ et $u^5 = uu' =
\alpha$ (pages 2 et 7) : $u$ engendre le corps $k[u] \simeq \mathbf{F}_{16}$,
il y est d'ordre multiplicatif $15$, et son image engendre le groupe cyclique
$k[u]^{*}/k^{*}$ d'ordre $5$. Le commutant de $k[u]$ dans $M$ est $k[u]$, et
$k[u] \cap S = \{u, u'\}$ ; le groupe d'ordre $5$ opère donc sans point fixe
sur les dix autres éléments de $S$, en deux orbites de cinq.

\emph{Traces des produits.} Pour $u, v \in S$, on a $\operatorname{Tr}(u^2) =
1$ et $\operatorname{Tr}(uu') = \operatorname{Tr}(\alpha\cdot 1) = 0$ ; et
$\operatorname{Tr}(uv') = \operatorname{Tr}(uv) + 1$. La formule de
polarisation donne
\[
  \det(u + v) = \det u + \det v + \operatorname{Tr}u\operatorname{Tr}v
  - \operatorname{Tr}(uv) = 1 + \operatorname{Tr}(uv)
\]
en caractéristique $2$\footnote{Les pages 2, 5 et 11 écrivent $\det(u+v) =
\det u + \det v + \operatorname{Tr}(uv)$, sans le terme
$\operatorname{Tr}u\operatorname{Tr}v = 1$, et la page 11 en tire
$\det(u+v) = \operatorname{Tr}(uv)$. Le tableau de la page 15 a la bonne
valeur. Les accolades « $1$ » sous $\det u$, $\det v$ (pages 2 et 5) sont aussi
de la page, en désaccord avec $\det u = \alpha$.}. La page 7 montre que, pour
$v$ fixé hors de $\{u, u'\}$, $\operatorname{Tr}(v\,u v u^{-1})$ ne prend
qu'une valeur par orbite du groupe d'ordre $5$, et pose la question de ces
valeurs.

\pagerange{9}{15}
\subsection*{Arêtes et coarêtes (pages 9 à 15)}

La page 9 choisit $u = \bigl(\begin{smallmatrix} 0 & 1 \\ \alpha & 1
\end{smallmatrix}\bigr)$ et un $v_0$ conjugué par une matrice diagonale ; la
page 11 traite les $v = u + \mu$ : $\det v = \alpha + \mu^2 + \mu$, qui vaut
$\alpha$ pour $\mu \in \{0, 1\}$ (c'est-à-dire $v \in \{u, u'\}$) et $\alpha'$
pour $\mu \in \{\alpha, \alpha'\}$. Le résultat est le tableau de la page 15 :
\[
  \begin{array}{c|cccc}
    & v = u & v = u' & \{u, v\} \text{ arête} & \{u, v\} \text{ coarête} \\ \hline
    \det(u - v) & 0 & 1 & \alpha' & \alpha \\
    \operatorname{Tr}(uv) & 1 & 0 & \alpha & \alpha'
  \end{array}
\]

\textbf{Proposition} (page 15). Soient $u, v \in S$, $v \notin \{u, u'\}$, et
$u^i(v) = u^i v u^{-i}$. Alors $\operatorname{Tr}(v\,u^i(v)) = \alpha$ si $i
\equiv \pm 1 \pmod 5$, $= \alpha'$ si $i \equiv \pm 2 \pmod 5$ (et $= 1$ si $i
\equiv 0$).

Autrement dit : la relation « $\operatorname{Tr}(uv) = \alpha$ » fait de $S$ le
graphe de l'icosaèdre, l'involution $u \mapsto u' = u + 1$ en est
l'antipodie\footnote{La page 5 écrit déjà « $u \mapsto u'$ antipodisme ».},
l'étoile d'un sommet $u$ est le pentagone $v_0, v_1 = u(v_0), \ldots, v_4$
parcouru par la rotation d'ordre $5$, et la relation
« $\operatorname{Tr}(uv) = \alpha'$ » relie chaque sommet aux cinq sommets à
distance $2$ : c'est le graphe d'un second icosaèdre sur les mêmes sommets
(le grand icosaèdre de Kepler--Poinsot)\footnote{Les noms « grand
icosaèdre », « Kepler--Poinsot » sont de nous ; la page dit
« coarête ». Le dénombrement de la page 15, $\binom{12}{2} = 66 = 6 + 30 +
30$, compte les paires antipodales, les arêtes et les coarêtes ; $\binom{12}{3}
= 220 = 20 + 20 + 60 + 60 + 60$ compte les triplets, les deux premiers $20$
étant les faces des deux icosaèdres.}. Pour un triangle $(u, v, w)$,
\[
  \det(u + v + w) = 3\alpha + 3\psi(u, v) = \alpha + \alpha' = 1,
\]
puisque $\psi(u, v) = 1 + \alpha = \alpha'$ sur une arête. La page 15 note enfin
que $(\alpha u)^5 = 1$ et $\det(\alpha u) = 1$ : $\alpha u$ est un relèvement
d'ordre $5$ de la rotation dans $\mathrm{SL}_2(\mathbf{F}_4)$.

\pagerange{13}{13}
\subsection*{Un calcul en caractéristique quelconque (page 13)}

Entre deux pages sur $\mathbf{F}_4$, la page 13 calcule avec des signes : pour
$u = \operatorname{diag}(\zeta, -\zeta)$ et $v$ de trace nulle, elle trouve
$uv + vu = \operatorname{Tr}(uv)\cdot 1$. C'est le cas $\operatorname{Tr}u =
\operatorname{Tr}v = 0$ de l'identité générale (6') de la page 81 (voir la
station IX).

\pagerange{18}{20}
\section*{II. La relation de tresses (pages 17 à 20)}

La page 17 est une chemise qui ne porte qu'au crayon « $\mathrm{sl}(2,
\Lambda)$ » ; les pages 18 à 20 sont à l'encre brune, datées « Nov. 1982 » et
« 19.6.83 ».

Soit $H$ un groupe engendré par deux éléments $\varepsilon_0$, $\varepsilon_1$.
On pose
\[
  \sigma = \varepsilon_0\varepsilon_1\varepsilon_0,\qquad
  \rho_0 = (\varepsilon_1\varepsilon_0)^{-1},\quad
  \rho_1 = (\varepsilon_0\varepsilon_1)^{-1},\qquad
  \rho_0' = \sigma\varepsilon_0,\quad \rho_1' = \sigma\varepsilon_1 .
\]
La page 20 vérifie que les relations $\sigma(\varepsilon_0) = \varepsilon_1$,
$\sigma(\varepsilon_1) = \varepsilon_0$ (conjugaison) et $\rho_0' =
\rho_0^{-2}$ équivalent chacune à la relation de tresses
$\varepsilon_0\varepsilon_1\varepsilon_0 = \varepsilon_1\varepsilon_0\varepsilon_1$,
tandis que $\rho_1' = \rho_1^{-2}$ est toujours vraie\footnote{La page écrit
$\rho_0^2$, $\rho_1^2$ là où les définitions donnent $\rho_0^{-2}$,
$\rho_1^{-2}$ ; le calcul s'arrête à cet endroit.}. Supposons donc
$H = B_3 = \langle \varepsilon_0, \varepsilon_1 \mid
\varepsilon_0\varepsilon_1\varepsilon_0 = \varepsilon_1\varepsilon_0\varepsilon_1
\rangle$, le groupe de tresses à trois brins\footnote{Nom de nous. La page 19
note ce groupe par un symbole lu $S\mathcal{T}^{+}_{1,1}$, de forme incertaine,
qui rappelle la notation $T_{g,\nu}$ des groupes de Teichmüller de
l'\emph{Esquisse d'un programme} ; $B_3$ est en effet le groupe des classes
d'isotopie d'un tore privé d'un disque. Le rapprochement est de nous.}. Alors :

\begin{itemize}
  \item la conjugaison par $\sigma$ échange $\varepsilon_0$ et
    $\varepsilon_1$, $\rho_0$ et $\rho_1$, $\rho_0'$ et $\rho_1'$ ; et
    $\rho_0(\varepsilon_0) = \rho_0'(\varepsilon_0) = \varepsilon_1$,
    $\varepsilon_0(\rho_0) = \rho_1$, etc. (points a) et a') de la page 18) ;
  \item $\sigma^2 = (\varepsilon_0\varepsilon_1)^3$ est central (la marge de
    la page 18 le note), $\rho_0^3 = \rho_1^3 = \sigma^{-2}$, et $\rho_0' =
    \rho_0^{-2} = \sigma^2\rho_0$ ;
  \item $B_3 \simeq \langle \rho, \sigma \mid \rho^{-3} = \sigma^2 \rangle
    \simeq \langle \rho, \varepsilon_0 \mid \rho\varepsilon_0\rho^{-1}\varepsilon_0\rho
    = 1\rangle$ (page 19, avec $\rho = \rho_0$) ;
  \item les listes h) à m) des pages 18 et 19 expriment les sept éléments en
    fonction de deux d'entre eux, et la note du 19 juin 1983 affirme que toute
    paire qui n'est pas formée de deux des quatre $\rho$ engendre $H$ ($21 - 6
    = 15$ paires), et que les quatre $\rho$ ne l'engendrent pas.
\end{itemize}

Les relations c) à e) de la page 18 — $\omega^2 = 1$ avec $\omega =
\sigma^{\pm 2}$, $\rho_0'^3 = 1$, $\rho_0^6 = 1$, $\rho_0' = \omega\rho_0$ — ne
valent pas dans $B_3$ mais dans son quotient par $\sigma^4$, et la marge le
dit (« valables dans $\mathrm{SL}(2, \mathbf{Z})$ mais non dans … »)\footnote{Le
second groupe de la marge est écrit comme le premier, $\mathrm{Sl}(2,
\mathbf{Z})$ ; le sens demande $B_3$. Les pages 18 et 20 posent l'une $\omega =
\sigma^{-2}$, l'autre $\omega = \sigma^2$ ; les deux coïncident dès que $\omega^2
= 1$.}. C'est le fait classique : $\varepsilon_0 \mapsto
\bigl(\begin{smallmatrix} 1 & 1 \\ 0 & 1\end{smallmatrix}\bigr)$,
$\varepsilon_1 \mapsto \bigl(\begin{smallmatrix} 1 & 0 \\ -1 & 1\end{smallmatrix}\bigr)$
définit un épimorphisme $B_3 \to \mathrm{SL}_2(\mathbf{Z})$ de noyau engendré
par $\sigma^4$, qui envoie $\sigma$ sur $\bigl(\begin{smallmatrix} 0 & 1 \\ -1
& 0\end{smallmatrix}\bigr)$ et $\omega = \sigma^2$ sur $-1$. La page 19 veut
en déduire que $\mathrm{GL}_2(\mathbf{Z})$ est une extension de
$\mathbf{Z}/2$ par $\mathrm{SL}_2(\mathbf{Z})$, à partir d'une extension de
$\mathbf{Z}/2$ par $B_3$ qu'« il faudrait expliciter »\footnote{La page dit
« par division par la relation $\omega = 1$ » ; diviser par $\omega$ donne
$\mathrm{PSL}_2(\mathbf{Z})$, diviser par $\omega^2$ donne
$\mathrm{SL}_2(\mathbf{Z})$.}.

Le paragraphe encadré en tête de la page 20 vient d'ailleurs : il parle d'un
sous-groupe commutatif $T_H$ de $H$, normalisant un $U_0$, avec $T_H \cap U_0
= \{1\}$ et $\sigma(t) = t^{-1}$, et dit que « la situation devient encore plus
jolie si on suppose $\zeta = 1$ ». C'est, mot pour mot, la situation
« particulièrement jolie » du § XI du dossier 79 (ses pages 238 à 247), où un
graphe homogène est décrit par un groupe $H$, un $U_0$, un $\sigma$ et un
$T_H$ ; voir la station suivante.

\pagerange{21}{28}
\section*{III. Brouillons dans $\mathrm{SL}_2$, et feuillets égarés (pages 21 à 28)}

\emph{Les calculs.} Les pages 21, 23, 25, 27 et le haut de la page 28 sont
des feuilles de calcul dans $H = \mathrm{SL}(M)$, $M = L \oplus L'$ somme de
deux modules inversibles sur $\Lambda$ ; elles préparent la rédaction des
pages 29 à 36, dont on reprend les notations à la station suivante : $e(\alpha)
= e_0(\alpha)$, $e'(\beta) = e_1(\beta)$ unipotents supérieur et inférieur,
$u(\lambda) = \operatorname{diag}(\lambda, \lambda^{-1})$, $N(c) =
\bigl(\begin{smallmatrix} 0 & -c^{-1} \\ c & 0\end{smallmatrix}\bigr)$. On y
trouve :

\begin{itemize}
  \item la grosse cellule (page 23) : $e(\alpha)u(\lambda)e'(\beta) =
    \bigl(\begin{smallmatrix} a & b \\ c & d\end{smallmatrix}\bigr)$ avec
    $\lambda = d^{-1}$, $\alpha = b/d$, $\beta = c/d$, donc $U_0 \times T
    \times U_1 \to H$ est injectif d'image $\{d \in \Lambda^{*}\}$ ;
  \item la formule d'échange (page 25) : si $1 + \alpha\beta \in \Lambda^{*}$,
    \[
      e'(\beta)e(\alpha) = e\Bigl(\frac{\alpha}{1+\alpha\beta}\Bigr)\,
      u\Bigl(\frac{1}{1+\alpha\beta}\Bigr)\,
      e'\Bigl(\frac{\beta}{1+\alpha\beta}\Bigr),
    \]
    et si $\alpha\beta = -1$, $e'(\beta)e(-\beta^{-1})e'(\beta) = N(\beta)$ ;
    la cellule $N^{-}B_1$ est $\{d = 0\}$ ;
  \item le produit de deux éléments de la grosse cellule (page 27) : si $1 +
    \alpha'\beta \in \Lambda^{*}$,
    \[
      e(\alpha)u(\lambda)e'(\beta)\cdot e(\alpha')u(\lambda')e'(\beta')
      = e\Bigl(\alpha + \frac{\lambda^2\alpha'}{1+\alpha'\beta}\Bigr)\,
      u\Bigl(\frac{\lambda\lambda'}{1+\alpha'\beta}\Bigr)\,
      e'\Bigl(\beta' + \frac{\lambda'^{2}\beta}{1+\alpha'\beta}\Bigr),
    \]
    et si $\alpha'\beta = -1$ le produit est $N(\beta\lambda^{-1}\lambda')\,
    e'(\beta' - \lambda'^2\beta - \alpha\beta^2\lambda'^2\lambda^{-2})$, grâce à
    $N(t)\,e(\alpha)\,N(t)^{-1} = e'(-t^2\alpha)$\footnote{La page écrit
    $\lambda'^{-2}$ aux deux endroits où l'on a ici $\lambda'^2$. On a
    $e'(y)u(\lambda') = u(\lambda')e'(\lambda'^2 y)$, comme le donne aussi la
    formule (14) de la page 31 ; c'est donc $\lambda'^2$.} ;
  \item le commutateur (page 23) : pour $\alpha\beta = 1$, $[e(\alpha),
    e'(\beta)] = \bigl(\begin{smallmatrix} 3 & -\alpha \\ \alpha^{-1} &
    0\end{smallmatrix}\bigr)$, qui ne tombe dans $N^{-}$ qu'en caractéristique
    $3$.
\end{itemize}

\emph{Les feuillets égarés.} Les pages 22 (tête-bêche), 24, 26 (numérotée
« 114 » de sa main) et le bas de la page 28 (numéroté « 135 ») appartiennent à
une autre suite, dont la numérotation des formules ((9) à (13), (64'), (70))
court hors du dossier. Elles parlent d'un groupe $G$ opérant sur un ensemble de
sommets $S$, d'un groupe $Z$ d'« homothéties », de sous-groupes $G_0^{+}
\subset G_0 \subset G$, d'un revêtement double $\widetilde{Z} = G/G_0^{+}$ de
$Z$, d'une arête orientée origine $(s_0, t_0)$, et du diagramme de
sous-groupes (13) :
\[
  \begin{gathered}
  B' = G_{s_0},\quad \widetilde{U} = B' \cap G_0,\quad U = B' \cap G_0^{+},\\
  \widetilde{B}'' = \{g \in G_0 \mid gs_0 \in Z s_0\},\quad
  B'' = \widetilde{B}'' \cap G_0^{+},\quad B = B'B'',
  \end{gathered}
\]
avec les indices $Z$, $\mu_2$ des inclusions successives. La page 22 commence
une « Remarque » sur un multigraphe $\Gamma = (S, \check{S}, R \subset S
\times \check{S})$ muni d'une opération libre d'un groupe commutatif $Z$ ;
la page 24 voudrait établir $\tilde{\delta}(\lambda) = \tilde{\lambda}^2$
pour $\tilde{\lambda} \in \widetilde{Z}$ d'image $\lambda$ (sa relation (70)) et
ne conclut pas ; la page 28 constate que $T_0$, $T_1$ et $T_G$ ont même image
dans $\operatorname{Aut}(H)$, qui « mérite d'être notée $\widetilde{Z}$ ». Ces
notations sont celles du § XI du dossier 79 (pages 222 à 269 de ce dossier :
multigraphes, homogénéité, le diagramme (64), $T_0$, $T_1$, $T_G$,
$T_H$)\footnote{Attribution de nous, sur la foi des notations et des numéros :
la pagination de l'auteur du dossier 79 va de 1 à 138, et ses pages 114 et
135 tombent dans son § XI. Nous n'avons pas vérifié que ces feuillets
comblent une lacune précise de ce dossier.}. On ne les restitue pas
davantage : elles supposent des définitions qui ne sont pas ici.

\pagerange{29}{36}
\section*{IV. Structure du groupe $\mathrm{SL}(2, \Lambda)$ (pages 29 à 36)}

Une rédaction suivie, paginée par lui de 1 à 4, formules (1) à (41) ; elle
s'arrête au milieu de la page 36.

\emph{Données} (pages 29 à 31). Soit $M$ un module libre de rang $2$ sur un
anneau $\Lambda$, muni d'une forme symplectique $e_M \in (\bigwedge^2 M)^{\vee}$,
de sorte que $\mathrm{SL}(M) = \operatorname{Aut}(M, e_M)$. Un tore maximal
déployé $T$ de $\mathrm{SL}(M)$ équivaut à une décomposition $M = L'_0 \oplus
L'_1$ en deux sous-modules de rang $1$ (ses deux droites propres), avec $L'_0
\otimes L'_1 \simeq \bigwedge^2 M \simeq \Lambda$ ; les caractères $\chi_i : T
\to \mathbb{G}_m$ par lesquels $T$ opère sur les $L'_i$ sont inverses l'un de
l'autre, et choisir l'un d'eux revient (quand $\operatorname{Spec}\Lambda$ est
connexe) à ordonner les deux droites\footnote{La marge de la page 29 nuance
le cas non connexe ; elle est en grande partie illisible.}. On écrit alors
\[
  \begin{gathered}
  H = \mathrm{SL}(M) = \Bigl\{\begin{pmatrix} a & b \\ c & d\end{pmatrix} \Bigm|
  a, d \in \Lambda,\ b \in L_0,\ c \in L_1,\ ad - bc = 1\Bigr\},\\
  L_0 = L_1'^{-1}L'_0,\qquad L_1 = L_0'^{-1}L'_1,
  \end{gathered}
\]
$L_0$ et $L_1$ étant inverses l'un de l'autre, $bc \in \Lambda$. On note
$u(\lambda) = \operatorname{diag}(\lambda, \lambda^{-1})$ ($\lambda \in
\Lambda^{*}$), $T = u(\Lambda^{*})$, et $e_0 : L_0 \to H$, $e_1 : L_1 \to H$ les
homomorphismes unipotents supérieur et inférieur, d'images $U_0$, $U_1$. Alors
$u(\lambda)e_0(\beta)u(\lambda)^{-1} = e_0(\lambda^2\beta)$ et
$u(\lambda)e_1(\gamma)u(\lambda)^{-1} = e_1(\lambda^{-2}\gamma)$\footnote{La
formule (15) écrit $e_0$ à la seconde ligne ; la marge souligne l'exposant
$2$.} ; $T \cap U_0 = T \cap U_1 = \{1\}$ ; les Borels $B_0 = TU_0 = \{c = 0\}$
et $B_1 = TU_1 = \{b = 0\}$ vérifient $B_0 \cap B_1 = T$\footnote{Dans (17), le
paramètre $\beta$ de $B_0$ est dit « $\in \Lambda$ » ; c'est $L_0$.}.

\emph{L'application $\sigma_0$} (pages 32 et 33). Supposons $L_0$ libre (sinon
$L_0^{*}$ est vide et ce qui suit est sans objet, dit la marge). On pose
\[
  \sigma_0 : L_0^{*} \to H,\qquad
  \sigma_0(\beta) = \begin{pmatrix} 0 & \beta \\ -\beta^{-1} & 0 \end{pmatrix},
\]
qui n'est \emph{pas} un homomorphisme. On a $\sigma_0(\beta)^2 = u(-1) = \omega$,
central ; $\sigma_0(\beta)$ normalise $T$ en y induisant $u(\lambda) \mapsto
u(\lambda^{-1})$ ; $\sigma_0(\lambda\beta) = u(\lambda)\sigma_0(\beta) =
\sigma_0(\beta)u(\lambda^{-1})$. Donc $N^{-} = \sigma_0(L_0^{*}) = \{a = d =
0\}$ est une classe à gauche et à droite modulo $T$, et $\sigma_0$ un
isomorphisme de $\Lambda^{*}$-torseurs. La « formule fondamentale » (25) est
\[
  \sigma_0(\beta)\,e_0(\beta)\,\sigma_0(\beta)^{-1} = e_1(-\beta^{-1}),
\]
d'où $\sigma(U_0) = U_1$ et $\sigma(U_1) = U_0$ pour tout $\sigma \in N^{-}$.

\emph{Transporteurs} (pages 33 et 34). La conjugaison d'un unipotent donne
\[
  g\,e_0(\beta)\,g^{-1} = \begin{pmatrix} 1 - ac\beta & a^2\beta \\ -c^2\beta & 1
  + ac\beta\end{pmatrix},\qquad g = \begin{pmatrix} a & b \\ c & d\end{pmatrix},
\]
d'où
\[
  \begin{gathered}
  \operatorname{Transp}_H(U_0, U_1) = \{a = 0\},\qquad
  \operatorname{Transp}_H(U_1, U_0) = \{d = 0\},\\
  N^{-} = \operatorname{Transp}(U_0, U_1) \cap \operatorname{Transp}(U_1, U_0),
  \end{gathered}
\]
et le transporteur est strict : si $a = 0$, $b$ et $c$ sont inversibles et
$\operatorname{int}(g)$ induit un isomorphisme $U_0 \simeq U_1$, linéaire
pour les structures de module\footnote{La marge de la page 33 ajoute qu'il
suffit que $g(U_0)$ rencontre $U_1$ en l'image d'\emph{un} $\beta \in L_0^{*}$
pour que $g$ soit dans le transporteur.}.

\emph{La relation de tresses} (pages 34 à 36). Pour $\beta \in L_0$, $\gamma \in
L_1$,
\[
  e_1(\gamma)\,e_0(\beta)\,e_1(\gamma) = \begin{pmatrix} 1 + \beta\gamma & \beta
  \\ \gamma(2 + \beta\gamma) & 1 + \beta\gamma\end{pmatrix},
\]
de sorte que les conditions suivantes sont équivalentes (37) : a)
$\beta\gamma = -1$ ; b) ce produit est dans $\operatorname{Transp}(U_0, U_1)$ ;
c) dans $\operatorname{Transp}(U_1, U_0)$ ; d) dans $N^{-}$ ; e) il est égal à
$\sigma_0(\beta)$. Pour $\beta\gamma = -1$, la formule (32)
$e_1(\gamma)e_0(\beta)e_1(\gamma) = \sigma_0(\beta)$ équivaut à
\[
  \rho(\beta)^3 = 1,\qquad \rho(\beta) = \sigma_0(\beta)\,e_0(\beta) =
  \begin{pmatrix} 0 & \beta \\ -\beta^{-1} & -1\end{pmatrix} = u(\beta)\,\rho(1),
\]
un élément de trace $-1$ et de déterminant $1$\footnote{L'élément
$e_1(\gamma)e_0(\beta)e_1(\gamma)$ avec $\gamma = -\beta^{-1}$ est l'élément
$w(t)$ de Steinberg pour la racine opposée ; la formule (32) et la relation
$\rho^3 = 1$ sont les relations de Steinberg de $\mathrm{SL}_2$. Noms de
nous.}. Il en sort une description de l'accouplement $L_0 \otimes L_1 \simeq
\Lambda$ en termes du seul groupe $H$ et de ses sous-groupes $U_0$, $U_1$
(page 35) : pour $u_0 = e_0(\beta) \in U_0^{*}$ et $u_1 = e_1(\gamma) \in U_1$,
\[
  \beta\gamma = -1 \iff u_1u_0u_1 \in N^{-} \iff u_0u_1u_0 = u_1u_0u_1 ,
\]
et l'on a alors $\sigma_0(u_0) = u_1u_0u_1$\footnote{La page 35 note que les
structures de module de $U_0$, $U_1$ sont automatiques quand $\Lambda$ est
$\mathbf{Z}$ ou un $\mathbf{Z}/n\mathbf{Z}$. L'équivalence avec la relation de
tresses demande $\beta$ inversible : la différence des deux produits a pour
coefficients non diagonaux $\pm\beta(1+\beta\gamma)$ et $\pm\gamma(1 +
\beta\gamma)$.}. La relation de tresses est donc le graphe d'une bijection
$u_0 \mapsto u_0^{\vee}$ de $U_0^{*}$ sur $U_1^{*}$, qui vérifie
$(u_0^{\lambda})^{\vee} = (u_0^{\vee})^{\lambda^{-1}}$ — alors que la
conjugaison par un élément du transporteur est linéaire. La rédaction
s'arrête sur la définition de $N' = T \sqcup N^{-}$, encadrée (41).

\pagerange{38}{50}
\section*{V. Rapport anharmonique et $n$-alités (pages 37 à 50)}

\pagerange{38}{40}
\subsection*{Trois flèches (pages 38 à 40)}

La couverture (page 37) porte « Rapport anharmonique — avril 1986 » ; la
rédaction est paginée par lui à partir de 1.

Soient $\mathcal{C}$ une catégorie abélienne, $I$ un ensemble à trois
éléments, et $\varphi_i : V \to V_i$ ($i \in I$) des flèches telles que pour
toute partie $J \subset I$ à deux éléments, $V \to \prod_{i \in J} V_i$ soit un
isomorphisme. Alors $0 \to V \to \prod_i V_i \to D \to 0$ est exacte, chaque
composé $u_i : V_i \to \prod V_i \to D$ est un isomorphisme, et le système est
isomorphe à $0 \to K(D, I) \to D^I \xrightarrow{\Sigma} D \to 0$, où $K(D, I)$
est l'ensemble des familles de somme nulle. La catégorie de ces systèmes
équivaut donc à celle des paires $(D, I)$ : $\mathrm{Ens}_3 \times \mathcal{C}$
(page 38)\footnote{Où $\mathrm{Ens}_3$ est le groupoïde des ensembles à trois
éléments. C'est la forme catégorique du fait qu'un « plan » muni de trois
« droites » deux à deux supplémentaires s'identifie à $D \oplus D$ avec les
droites $D \times 0$, $0 \times D$ et la diagonale.}.

Il y a une variante non commutative (pages 38 à 40) : les $V_i$ sont des
ensembles, chaque $V_i$ porte une structure de bitorseur sous deux des trois
groupes $G_j$, définis à isomorphisme unique près, et $V$ s'identifie aux
triplets $(x_i, x_j, x_h)$ tels que $x_i \wedge x_j \wedge x_h = \varepsilon$,
pour un élément central $\varepsilon$ du produit contracté $V_i
\wedge_{G_k} V_j \wedge_{G_i} V_h$, une fois choisie une orientation de $I$.
Si l'un des $G_i$ est commutatif, tous le sont et l'on revient au cas
précédent dès qu'on a choisi une origine. Le groupe des automorphismes se
dévisse en $\mathfrak{S}_I \simeq \mathfrak{S}_3$, en
$\operatorname{Aut}\operatorname{Ext}(G_i)$ et en $G^I/Z$ (marge de la page
40)\footnote{Lecture mince : la page 40 est rapide et la prose de liaison y est
incertaine.}.

\pagerange{40}{45}
\subsection*{Quatre flèches et leurs rapports anharmoniques (pages 40 à 45)}

Soit maintenant $I$ à quatre éléments, avec la même condition sur les paires.
Pour toute numérotation $I = \{i_1, i_2, i_3, i_4\}$, l'isomorphisme
canonique des trois flèches $i_1, i_2, i_3$ (en oubliant $\varphi_{i_4}$)
donne $u_{(i_1 i_2)(i_3 i_4)} : V_{i_1} \simeq V_{i_2}$, avec les relations (2)
et (3) de la page 41 ; $V$ se reconstruit comme sous-objet de $\prod V_i$ à
partir du système des $u$, à des conditions près que la page 41 annonce. On
définit (page 42) le rapport anharmonique
\[
  \rho_{ijkl} = u_{ijlk}^{-1}\,u_{ijkl} \in \operatorname{Aut}(V_i),
\]
avec $\rho_{ijlk} = \rho_{ijkl}^{-1}$, $\rho_{jilk} = u_{ijkl}(\rho_{ijkl})$
(« invariance par double permutation »), et la relation
\[
  (7)\qquad \rho_{ikjl} = 1 - \rho_{ijkl},
\]
que la marge demande de vérifier\footnote{Nous ne l'avons pas vérifiée dans
la généralité de la page, où le signe dépend des conventions de définition des
$u$ ; pour quatre droites deux à deux distinctes d'un plan, c'est l'une des
relations classiques du birapport, à une convention près.}. La proposition
des pages 43 et 44 énonce qu'un système de $u$ vérifiant (2), (3)
provient d'une « quadralité » $(V, (V_i), (\varphi_i))$ si et seulement si les
vingt-quatre relations (7) ont lieu ; elle n'est pas démontrée. Quand les
$\operatorname{End}(V_i)$ sont commutatifs, ils s'identifient canoniquement à
un même anneau $k$, et les relations deviennent
\[
  \rho_{jikl} = \rho_{ijlk} = \rho_{ijkl}^{-1},\qquad \rho_{ikjl} = 1 - \rho_{ijkl}
  \qquad(\Rightarrow\ \rho_{jilk} = \rho_{ijkl}) :
\]
$(i, j, k, l) \mapsto \rho_{ijkl}$ est une application
$\mathfrak{S}_4$-équivariante du torseur des numérotations de $I$ dans
$U(k) = \{\lambda \mid \lambda,\ 1 - \lambda \in k^{*}\}$, où $\mathfrak{S}_4$
opère par son quotient $\mathfrak{S}_3$ ($\lambda \mapsto \lambda^{-1}$, $1 -
\lambda$). Autrement dit (page 45), $\rho$ ne dépend que d'un « coin » — une
partition de $I$ en deux paires munie d'un ordre — et la donnée de $\rho$
équivaut à une application du torseur des repères de l'ensemble à trois
éléments des partitions de $I$ dans $U(k)$\footnote{C'est le noyau de Klein
$\mathfrak{V}_4 \subset \mathfrak{S}_4$ qui agit trivialement : la
description classique du birapport. La page conclut qu'il « reste à décrire
la catégorie des quadralités … de façon plus simple », et ne le fait pas.}.

\pagerange{47}{50}
\subsection*{$n$-alités ensemblistes (pages 47 à 50)}

Datée « Avril 86 ». Une \emph{$n$-alité} dans une catégorie $C$ est un système
$(I, E, (E_i)_{i \in I}, (\varphi_i : E \to E_i))$, $\operatorname{card} I =
n$, tel que pour toute paire $J \subset I$, $E$ muni des $\varphi_i$ ($i \in
J$) soit un produit des $E_i$ ; elles forment une catégorie fibrée sur le
groupoïde des ensembles à $n$ éléments. Pour $C = \mathrm{Ens}$, hors du cas
« impropre » (un $E_i$ vide) et du cas « trivial » (tous réduits à un point),
c'est un ensemble $E$ muni de $n$ partitions dont deux quelconques sont
« orthogonales »\footnote{Pour $n = 3$ c'est un carré latin (un quasi-groupe à
isomorphisme près), pour $n$ quelconque un réseau (\emph{net}) d'ordre
$\operatorname{card}E_i$, c'est-à-dire $n - 2$ carrés latins deux à deux
orthogonaux. Noms de nous.}. La page 48 le traduit en structure d'incidence
entre les « sommets » $S = \coprod E_i$ et les points de $E$, $R \subset S
\times E$, et dégage les axiomes : (i) deux sommets distincts sont incidents à
au plus un $x$ (condition symétrique en $S$ et $E$) ; (ii) « ne pas être une
arête » est une relation d'équivalence sur $S$, dont les classes sont les
$E_i$ ; (iii) chaque $R(x)$ rencontre chaque classe ; (iv) un $R(x)$ au moins a
deux éléments\footnote{C'est l'axiomatique d'un plan transversal
(\emph{transversal design}) $\mathrm{TD}(n, m)$, dual d'un réseau. Nom de nous.}.

\pagerange{52}{60}
\section*{VI. Bi-icosaèdres et cartes de congruence (pages 51 à 60)}

La couverture (page 51) porte « Bi-icosaèdres et cartes de congruences, avril
1986 ». La lecture des pages 52 à 60 est rapide ; on n'en restitue que ce que
la transcription porte sûrement.

\pagerange{52}{55}
\subsection*{Orientations (pages 52 à 55)}

Une structure de \emph{bi-icosaèdre} sur un ensemble $S$ de douze sommets est
une paire $\{\mathcal{I}, \mathcal{I}'\}$ de structures d'icosaèdre
mutuellement complémentaires : $\{s, s'\}$ est une arête de $\mathcal{I}'$ si et
seulement si $\{s, -s'\}$ en est une de $\mathcal{I}$, où $s \mapsto -s$ est
l'antipodie (page 55). Autour d'un sommet $s$, le groupe des rotations $U_s$ est
un $\mathbf{F}_5$-module libre de rang $1$ ; ses générateurs forment un torseur
sous $\mathbf{F}_5^{*} \simeq \mathbf{Z}/4$, dont les deux classes modulo $\pm
1$, $\{u, -u\}$ et $\{2u, -2u\}$, sont les deux orientations de l'étoile de $s$
pour $\mathcal{I}$ et pour $\mathcal{I}'$ (page 52). L'application $u \mapsto
2u$ identifie les revêtements d'orientation des deux icosaèdres au-dessus des
sommets, mais l'aller-retour vaut $u \mapsto 4u = -u$, l'antipodie : sur les
sommets, l'identification dépend du choix de $\mathcal{I}$ ou de $\mathcal{I}'$.
Sur les faces il n'y a pas d'ambiguïté (une face de $\mathcal{I}$ et sa
« complémentaire » se correspondent canoniquement, l'aller-retour est
l'identité) ; sur les arêtes, une bijection entre arêtes orientées revient au
choix d'une orientation de l'arête commune (page 53). D'où la conclusion de la
page 53 : les dodécaèdres duaux $D$, $D'$ s'identifient canoniquement, les
icosaèdres $I$, $I'$ ne le font qu'au produit près par l'antipodie. Un
bi-icosaèdre \emph{orienté} est défini (page 54) comme une paire d'icosaèdres
orientés avec une paire $\{\Theta, \Theta'\}$ de bijections de sommets
compatibles aux antipodies, différant par l'antipodie, ou encore par des
isomorphismes $\vec{I}' \simeq \operatorname{compl}(\vec{I})$, $\vec{I}
\simeq \operatorname{compl}(\vec{I}')$ dont le composé est l'antipodie.

\pagerange{55}{59}
\subsection*{Les groupes (pages 55 à 59)}

Le groupe des automorphismes d'un bi-icosaèdre s'envoie sur le groupe des
permutations de $\{\mathcal{I}, \mathcal{I}'\}$ :
\[
  1 \to \operatorname{Aut}(\mathcal{I}) \to \operatorname{Aut}(\mathcal{I},
  \mathcal{I}') \to \mathbf{Z}/2 \to 1,\qquad
  \operatorname{Aut}(\mathcal{I}) \simeq \operatorname{Aut}^{+}(\mathcal{I})
  \times \{\pm 1\},
\]
et le quotient de $\operatorname{Aut}(\mathcal{I}, \mathcal{I}')$ par le
sous-groupe qui respecte $\mathcal{I}$, $\mathcal{I}'$ et l'orientation est
cyclique d'ordre $4$, « sauf erreur » (page 56). La page 57 compare deux
extensions de noyau $\operatorname{Aut}(\vec{I})$, l'une par $\mathbf{Z}/2$
(bi-icosaèdres orientés), l'autre par $\mathbf{Z}/4$, et constate qu'il n'y a
pas de flèche de l'une dans l'autre induisant l'identité sur le noyau. La page 58
réalise tout dans la théorie des cartes de congruence :
\[
  1 \to \mathrm{SL}_2(\mathbf{F}_5)/\{\pm 1\} \to \mathrm{GL}_2(\mathbf{F}_5)/\{\pm 1\}
  \xrightarrow{\det} \mathbf{F}_5^{*} \simeq \mathbf{Z}/4 \to 1 .
\]
Ici $S$ est l'ensemble des vecteurs primitifs de $\mathbf{F}_5^2$ modulo $\pm
1$ ; les arêtes de $\mathcal{I}$ sont les $\{\dot{a}, \dot{b}\}$ avec $\det(a, b)
= \pm 1$, celles de $\mathcal{I}'$ les $\{\dot a, \dot b\}$ avec $\det(a, b) =
\pm 2$, et l'antipodie est l'homothétie de rapport $2$\footnote{Explicitation
de nous, conforme au dossier 79 (ses pages 37 à 60, « deux icosaèdres
entrelacés »). Elle confirme le « sauf erreur » : $\operatorname{Aut}(\mathcal{I},
\mathcal{I}') \simeq \mathrm{GL}_2(\mathbf{F}_5)/\{\pm 1\}$, d'ordre $240$, et son
quotient par $\mathrm{SL}_2(\mathbf{F}_5)/\{\pm 1\} \simeq \mathfrak{A}_5$ est
$\mathbf{F}_5^{*} \simeq \mathbf{Z}/4$. Le quotient par l'antipodie (la classe
de l'homothétie $2$) est $\mathrm{PGL}_2(\mathbf{F}_5) \simeq \mathfrak{S}_5$,
d'ordre $120$ : c'est sans doute ce que vise la page 58 quand elle écrit
« $\mathrm{PGl}(2, \mathbf{F}_5) \times \mathbf{F}_5^{*}$ ($\simeq
\mathfrak{S}_5$, ordre $120$) », formule qui, telle quelle, ne tient pas.}.
Le quotient d'ordre $8$ de la page 58 opère sur les quatre orientations de
$(\mathcal{I}, \mathcal{I}')$ comme le groupe des automorphismes d'un carré, et
la page 59 lit les deux extensions comme les images réciproques de deux
sous-groupes de ce groupe diédral.

\pagerange{59}{60}
\subsection*{Généralisation : le cône nilpotent (pages 59 et 60)}

Pour $n$ quelconque, la catégorie $\mathrm{Congr}$ des cartes de congruence se
réalise sur les paires $\{a, -a\}$ de vecteurs primitifs d'un
$\mathbf{Z}/n$-module symplectique $E$ libre de rang $2$ ; on envoie $a$ sur
$a \otimes a \in \operatorname{End}(E) = M_2(\mathbf{Z}/n)$, l'endomorphisme
$x \mapsto (a \wedge x)\,a$ (via l'identification symplectique $E \simeq
\check{E}$), qui est de trace et de carré nuls. On réalise ainsi $S$ dans le
cône des nilpotents, de façon compatible à $\mathrm{SL}(E)/\{\pm 1\}$ mais non
à $\mathrm{GL}(E)$ ; pour $n = p$ premier on n'obtient pas tous les nilpotents
non nuls, mais une classe sur deux modulo $\mathbf{F}_p^{*}/\mathbf{F}_p^{*2}$,
l'homothétie par $\lambda$ sur les vecteurs devenant l'homothétie par
$\lambda^2$ sur les nilpotents (page 60). La page note que cette division en
classes a un sens pour toute algèbre d'Azumaya de rang $4$ sur un corps, et
s'interrompt sur « TSVP ».

\pagerange{61}{73}
\section*{VII. Le cône nilpotent d'une algèbre d'Azumaya (pages 61 à 73)}

Les pages 61 à 73 sont très raturées et chargées de notes obliques en grande
partie illisibles ; la lecture qui suit s'appuie sur les formules et les
énoncés que la transcription donne sûrement, et elle est mince aux pages 61
à 63 et 68.

\pagerange{61}{63}
\subsection*{Classes de nilpotents et carte universelle (pages 61 à 63)}

Soit $A$ une algèbre d'Azumaya de rang $4$ sur $\mathbf{F}_p$, $p$ impair, et
$C^{*}$ l'ensemble de ses nilpotents non nuls. Il se divise en deux classes
$C^{*}_i$ avec $\lambda C^{*}_i = C^{*}_i$ si et seulement si $\lambda \in
k^{*2}$ ; chacune porte une structure canonique de carte orientée, et
$A^{*}/k^{*}$ opère en respectant la structure, le sous-groupe d'indice $2$
préservant chaque classe (page 61). Pour $p = 5$, $\mathbf{F}_5^{*2} = \{\pm 1\}$
et l'on retrouve la structure de bi-icosaèdre (page 62)\footnote{La page 61
met en garde : $\overline{T}_\lambda$ \emph{transporte} les structures, $\{a,
b\}$ est une arête de $T_\lambda(c_i)$ si $\{T_\lambda a, T_\lambda b\}$ en est
une de $c_i$, « et non pas » si $\{a, T_\lambda b\}$ en est une. Une marge
note une classe « bien mystérieuse ».}.

La « carte universelle », définie sur $\mathbf{Z}$, est réalisée dans l'algèbre
engendrée par un $\sigma$ de déterminant $1$ et de trace $0$ (d'ordre $4$) et un
$\rho_f$ de déterminant $1$ et de trace $-1$ (d'ordre $3$), avec $\sigma\rho_f
= \rho_s^{-1}$ unipotent ; on obtient une représentation fidèle de
$\mathrm{SL}_2(\mathbf{Z})/\{\pm 1\}$\footnote{La page écrit aussi « $\rho_f$
d'ordre $6$ », ce qui est l'ordre de $-\rho_f$ ; dans $\mathrm{SL}_2$, un
élément de trace $-1$ est d'ordre $3$. Selon le choix des signes, le produit
$\sigma\rho_f$ est unipotent ou opposé d'un unipotent ; modulo $\pm 1$, c'est
indifférent. Le groupe engendré est le groupe cartographique triangulé du
dossier 79 (ses pages 26 à 37).}, qui se factorise modulo $n$ par
$\mathrm{SL}_2(\mathbf{Z}/n)/\{\pm 1\}$ et redonne la réalisation précédente
(page 63). La marge de la page 63 en donne le dictionnaire : les sommets sont
les unipotents, les arêtes les paires $(u, u')$ d'unipotents avec
$\operatorname{Tr}(uu') = -1$\footnote{Pour des unipotents $1 + n$, $1 + n'$ ($n$,
$n'$ nilpotents), $\operatorname{Tr}((1+n)(1+n')) = 2 + \operatorname{Tr}(nn')$ :
la condition se lit sur les nilpotents comme $\operatorname{Tr}(nn') = -3$ ;
c'est la forme $\operatorname{Tr}(nn') = -1$ des pages 68 et 73 que la suite
retient, pour les nilpotents eux-mêmes. Remarque de nous.}.

\pagerange{64}{68}
\subsection*{Les classes $C^{*}(e)$ (pages 64 à 68)}

Soit maintenant $k$ un anneau et $A \simeq \operatorname{End}(E)$, $E$
projectif de rang $2$, $\underline{\omega} = \det E$. Un élément $u$ de carré nul
non dégénéré définit une droite $D = ku \simeq L^{\otimes 2} \otimes
\underline{\omega}^{-1}$, où $L$ est son image\footnote{La page passe par
$L' = L^{\perp} \subset \check{E} \simeq E \otimes \underline{\omega}^{-1}$ et
$L' \simeq L \otimes \underline{\omega}^{-1}$, « une fois faite la convention de
signe ».}. Le choix d'une base $e$ de $\underline{\omega}$ (une structure
symplectique sur $E$) donne $\varphi^{e} : L^{\otimes 2} \simeq D$, avec
$\varphi^{\lambda e} = \lambda\varphi^{e}$ ; si $L$ est libre, $D^{*}$ contient
une classe privilégiée modulo $k^{*2}$, $D^{*}(e) = \{\varphi^{e}(x \otimes x)
\mid x \in L^{*}\}$, et $D^{*}(\lambda e) = \lambda D^{*}(e)$. Explicitement,
$\varphi^{e}(x \otimes x)$ est l'endomorphisme
\[
  x \overset{e}{\otimes} x : y \longmapsto \frac{x \wedge y}{e}\,x ,
\]
de trace et de carré nuls. Sous l'hypothèse ${}_2\mathrm{Pic}(k) = 0$, toutes
les droites $L_u$ sont libres, et l'on pose $C^{*}(e) = \bigcup_D D^{*}(e)$ ;
alors
\[
  C^{*} = \coprod_{e \in \underline{\omega}^{*}/k^{*2}} C^{*}(e),
\]
partition indexée par un torseur sous $k^{*}/k^{*2}$, permutée transitivement
par les homothéties, et qui ne dépend que de $A$ (page 66). L'application
$\varphi^{e} : E^{*} \to C^{*}(e)$, $x \mapsto x \overset{e}{\otimes} x$, est
surjective, avec $\varphi^{e}(\lambda x) = \lambda^2\varphi^{e}(x)$, et pour
$\Theta \in \mathrm{GL}(E)$ opérant par conjugaison,
\[
  \Theta\,\varphi^{e}(x)\,\Theta^{-1} = (\det\Theta)\,\varphi^{e}(\Theta x)
\]
(page 67, encadré) : la compatibilité n'a lieu que sur $\mathrm{SL}(E)$.
Chaque classe porte une structure triangulaire orientée canonique,
indépendante de $e$, et une « $\lambda$-orientation » qu'« il faudrait
expliciter ». Les remarques de la page 68 : $C^{*}$ n'est non vide que si $A
\simeq \operatorname{End}(k \oplus L)$ ; et il existe une arête (deux éléments
de $C^{*}$ avec $\operatorname{Tr}(uv) = -1$) si et seulement si $A \simeq
M_2(k)$, car un tel couple est conjugué au couple type de $M_2$.

\pagerange{69}{73}
\subsection*{Triangles et orientation (pages 69 à 73)}

Soient $x, y, z \in E$ avec $x + y + z = 0$ et $x \wedge y = y \wedge z = z \wedge
x = e$. Il existe un unique $r \in \mathrm{SL}(E)$ tel que $rx = y$, $ry = z$, $rz
= x$ ; il est d'ordre $3$, ne dépend que de $\{\dot x, \dot y, \dot z\}$ et d'un
ordre circulaire, et changer l'ordre change $r$ en $r^{-1}$. On a
$\det r = 1$ et $\operatorname{Tr}r = -1$, donc $r^{-1} = r' = -1 - r$\footnote{La
page écrit $\operatorname{Tr}\rho = 1$, $\rho^{-1} = 1 - \rho$ et $\det(1 -
2\rho) = 3$ : ce sont les formules de $-r$, d'ordre $6$, alors qu'elle dit
« d'ordre $3$ ». La transcription le signale. On rétablit $r$ d'ordre $3$.}.
L'élément $\tau = x \overset{e}{\otimes} x + y \overset{e}{\otimes} y + z
\overset{e}{\otimes} z$ commute à $r$, a une trace nulle, et
\[
  \det\tau = 3\cdot 0 + \psi(u, v) + \psi(v, w) + \psi(w, u) = 3
\]
puisque $\psi = -\operatorname{Tr}(uv) = 1$ sur chaque arête ; donc $\tau = a r
+ b$ avec $a = 2b$, $\det\tau = 3b^2$, $b^2 = 1$ et
\[
  \tau = \pm(1 + 2r)
\]
(pages 69 et 70)\footnote{Dans les conventions de la page, $\pm(1 - 2\rho)$
avec $\rho = -r$. Le signe effectif dépend de la convention qui identifie $x
\otimes x$ à un endomorphisme ; avec celle qu'on a écrite plus haut, le
calcul direct donne $\tau x = z - y = -(1 + 2r)x$.}. Comme $1 + 2r^{-1} = -(1 +
2r)$, quand $2$ n'est pas diviseur de zéro le signe fixé choisit l'un de $r$,
$r^{-1}$ : il y a une seule rotation $r$ telle que $u + v + w = \varepsilon(1 +
2r)$ pour le signe $\varepsilon$ de la convention, et elle définit l'orientation
circulaire canonique du triangle $\{u, v, w\}$ de $C^{*}(e)$. La page ne
tranche pas « lequel ! ».

Si $\{\pm 1\} \simeq \mu_2(k)$ (et ${}_2\mathrm{Pic}(k) = 0$), on a
$\varphi^{e}(x) = \varphi^{e}(y)$ si et seulement si $y = \pm x$, et
$\psi^{e} : E^{*}/\{\pm 1\} \to C^{*}(e)$ est bijectif (page 71) ; il
transporte sur $C^{*}(e)$ la structure de bi-polyèdre orienté du dossier 79 :
$\{\dot x, \dot y\}$ est une arête si $x \wedge y = \pm e$, les triangles sont
les $\{\dot x, \dot y, \dot z\}$ avec $x + y + z = 0$, chaque arête est dans
deux triangles qui l'orientent en sens opposés, et la translation $T_\lambda(\dot
y) = \overline{y + \lambda x}$ de l'étoile de $\dot x$ est induite par la
transvection $1 + \lambda\varphi^{e}(x)$ (pages 72 et 73)\footnote{La
condition $\{\pm 1\} \simeq \mu_2(k)$ équivaut, si $2$ est inversible, à la
connexité de $\operatorname{Spec}k$ ; le crochet de la page 71 traite le cas
d'un nombre fini de composantes. Le diagramme de la page 71 identifie
$E^{*}/k^{*} \simeq \mathbb{P}(E)(k)$ à $C^{*}/k^{*}$, « bien connu ».}. Dans
$A$ :
\[
  (x \overset{e}{\otimes} x)(y \overset{e}{\otimes} y) = \frac{x \wedge y}{e}\,(x
  \overset{e}{\otimes} y),\qquad
  \operatorname{Tr}(uv) = -\Bigl(\frac{x \wedge y}{e}\Bigr)^2,
\]
de sorte que $x \wedge y = \pm e \iff \operatorname{Tr}(uv) = -1$ : les arêtes de
$C^{*}(e)$ sont les paires de traces de produit $-1$, les triangles les
triplets deux à deux tels\footnote{La page écrit le composé $y
\overset{e}{\otimes} x$ ; la différence est de convention, la trace est la même.
La dernière marge de la page 73 demande ce que donnerait l'autre relation,
$\operatorname{Tr}(uv) = +1$.}. Le fil s'interrompt là.

\pagerange{75}{77}
\section*{VIII. Tentatives : $A$ décrite par $A^{*}/k^{*}$ (pages 74 à 77)}

La page 74 est une couverture, « Azumaya rang 4 — avril 1986 ». Les pages 75
et 77 cherchent, titre de la page 77, une « description élémentaire de $A \simeq
M_2(k)$ ($k$ corps) en termes de $G = A^{*}/k^{*}$ ». Sur $G$, la fonction
\[
  \delta(\bar x) = \frac{b^2 - 4c}{c},\qquad b = \operatorname{Tr}x,\ c = \det x,
\]
est bien définie (invariante par $x \mapsto \lambda x$), centrale
($\delta(xy) = \delta(yx)$), paire ($\delta(x^{-1}) = \delta(x)$), nulle en $1$,
et $\delta(x^2) = \delta(x)^2 + 4\delta(x)$. La page propose
\[
  \delta(xy) + \delta(xy^{-1}) = \delta(x)\delta(y) + 2(\delta(x) + \delta(y))
\]
avec en marge « si $x, y$ commutent seulement ??? — À vérifier ! ». La
relation exacte est
\[
  \delta(xy) + \delta(xy^{-1}) = \delta(x)\delta(y) + 2(\delta(x) + \delta(y)) +
  2(\kappa - 2),\qquad \kappa = \operatorname{Tr}(xyx^{-1}y^{-1}),
\]
qui redonne celle de la page exactement quand le commutateur a pour trace
$2$ — par exemple quand $x$ et $y$ commutent, ou ont un vecteur propre
commun\footnote{Notre calcul : avec $s(x) = \delta(x) + 4 =
(\operatorname{Tr}x)^2/\det x$, c'est l'identité de Fricke
$\operatorname{Tr}(xy)\operatorname{Tr}(xy^{-1}) = (\operatorname{Tr}x)^2 +
(\operatorname{Tr}y)^2 - 2 - \kappa$ dans $\mathrm{SL}_2$, élevée au carré. Vérifié
aussi numériquement.}. La page demande aussi un polynôme $P \in
\mathbf{Z}[X, Y, Z]$ tel que $[x, y]^2 - P(\delta x, \delta y, \delta(xy))[x, y] +
1 = 0$, c'est-à-dire $\kappa = P(\ldots)$ ; il n'existe pas : $\kappa$ fait
intervenir $\operatorname{Tr}x\operatorname{Tr}y\operatorname{Tr}(xy)/(\det x\det
y)$, dont seul le carré, $(\delta x + 4)(\delta y + 4)(\delta(xy) + 4)$, est
polynomial en les $\delta$\footnote{Remarque de nous. De même, la relation
cherchée à la page 77 entre les dix quantités $\delta(x_i)$, $\delta(x_ix_j)$
pour quatre éléments, et le polynôme $Q$ tel que $\delta(xyz) =
Q(\delta x, \delta y, \delta z, \delta(xy), \delta(yz), \delta(zx))$, ne peuvent
exister sous cette forme : la trace de $xyz$ n'est déterminée par les traces
plus courtes qu'au choix près d'une racine d'une équation du second degré.
C'est exactement ce que la station XII établira.}. La page 75 tire de trois
parenthésages de $\delta(xyzt)$ deux relations entre dix quantités, puis se
demande si $\delta(\tau h) = \delta(\tau)$ pour $\tau$ normalisant un sous-groupe
commutatif $H$ ; la marge répond « C'est idiot ! » (prendre $\tau \in H$).

\pagerange{80}{99}
\section*{IX. Azumaya rang 4 (I) : l'algèbre $M(b, c; b', c'; t)$ (pages 80 à 99)}

\pagerange{80}{84}
\subsection*{Hamilton--Cayley et l'algèbre engendrée par deux éléments (pages 80 à 84)}

Soit $A$ d'Azumaya de rang $4$ sur un anneau $k$. On a $\operatorname{Tr}$
linéaire, $\det$ multiplicatif et quadratique, Hamilton--Cayley
$u^2 - \operatorname{Tr}(u)u + \det(u) = 0$\footnote{La page 80 écrit « $=
1$ » au second membre.}, la polarisation
\[
  (3)\qquad \det(u + v) = \det u + \det v + \operatorname{Tr}u\operatorname{Tr}v -
  \operatorname{Tr}(uv),
\]
et $\operatorname{Tr}1 = 2$, $\det 1 = 1$. La remarque de la page 80 : pour une
algèbre projective de type fini où $2$ est inversible, (1) et (2) entraînent
(3), via $(\operatorname{Tr}u)^2 - \operatorname{Tr}(u^2) = 2\det u$. En
polarisant Hamilton--Cayley :
\[
  (6')\qquad uv + vu = \operatorname{Tr}(v)\,u + \operatorname{Tr}(u)\,v +
  \bigl(\operatorname{Tr}(uv) - \operatorname{Tr}u\operatorname{Tr}v\bigr)\,1,
\]
en particulier $vu \equiv -uv$ modulo $k1 + ku + kv$\footnote{La page 81 écrit
$vu \equiv uv$ ; la transcription relève l'écart.}.

\textbf{Proposition} (pages 81 et 82). Dans une algèbre vérifiant (1), (2),
(3), le sous-module engendré par $1, u, v, uv$ est une sous-algèbre. Il suffit
de vérifier la stabilité par multiplication à gauche par $u$ et $v$, ce que
donnent Hamilton--Cayley et (6').

On en tire (page 83) une algèbre $M(b, c; b', c'; t)$ libre de base $1, u, v,
uv$ pour tout $(b, c, b', c', t) \in k^5$, définie sur
$\mathbf{Z}[B, C, B', C', T]$, de table de multiplication (page 92 ; facteur de
gauche en ligne)
\[
  \begin{array}{c|ccc}
    & u & v & uv \\ \hline
    u & bu - c & uv & b\,uv - cv \\
    v & (t - bb') + b'u + bv - uv & b'v - c' & -bc' + c'u + tv \\
    uv & -b'c + tu + cv & b'uv - c'u & t\,uv - cc'
  \end{array}
\]
L'associativité, la multiplicativité de $\det$ défini par (10) sur $k[u]
\oplus k[v]$, et les relations (2), (3) se « vérifient » ; la page 84 évite le
calcul par prolongement des identités algébriques, en se plaçant sur l'ouvert
où $\delta$ et $\delta'$ sont inversibles et en réalisant presque tout $t$
dans $M_2$\footnote{Nous avons vérifié la table ; l'associativité aussi se
vérifie directement. La marge de la page 122 met en garde contre un cercle
vicieux à éviter si l'on veut se dispenser de cette vérification.}.

\pagerange{85}{91}
\subsection*{L'invariant anharmonique (pages 85 à 91)}

\emph{Lemme} (page 85). Si $u$, $v \in A$ engendrent des sous-algèbres
quadratiques étales, $\Theta \mapsto \operatorname{Tr}(u\,\Theta v\Theta^{-1})$
est un morphisme dominant vers la droite affine.

Deux sous-algèbres étales $A_1 = k[u]$, $A_2 = k[v]$ définissent chacune un
diviseur de degré $2$ sur la conique associée à $A$ ; si leurs supports sont
disjoints, ces deux paires de points ont un rapport anharmonique $\lambda$,
défini localement et au remplacement près de $\lambda$ par $\lambda^{-1}$, d'où
l'invariant $\rho = \lambda + \lambda^{-1}$, défini sur $k$, racine de
$\lambda^2 - \rho\lambda + 1 = 0$ (page 86). La page 87 le calcule :
\[
  (11)\qquad \rho = 2\,\frac{(2t - bb')^2 + \delta\delta'}{(2t - bb')^2 -
  \delta\delta'},\qquad
  \rho_0 = \rho - 2 = \frac{\delta\delta'}{Q},
\]
avec $Q = t^2 - bb't + b^2c' + b'^2c - 4cc'$, puisque $(2t - bb')^2 -
\delta\delta' = 4Q$ ; en caractéristique $2$, $\rho = (bb')^2/(t^2 + bb't + b^2c'
+ b'^2c)$\footnote{Formules vérifiées. Pour $b = b' = 0$ la page trouve $\rho =
2(t^2 + 4cc')/(t^2 - 4cc')$, et $t^2 - 4cc'$ est alors le discriminant de $uv$.
La page 86 écrit la condition « $P - 2Q \neq 2$ partout » ; on lit $P - 2Q =
\delta\delta'$ partout inversible.}. Les valeurs exclues se lisent sur
$\lambda$ : $\lambda = 1$ (soit $\rho = 2$) correspond à $\delta\delta' = 0$, et
$\lambda \in \{0, \infty\}$ (soit $\rho = \infty$) à $Q = 0$, une paire de points
fixes commune\footnote{Avec la convention $\lambda = (q_1 - p_1)(q_2 -
p_2)/((q_1 - p_2)(q_2 - p_1))$ ; explicitation de nous.}.

\textbf{Proposition} (page 89). Si $\delta\delta'Q$ est inversible, $M(b, c;
b', c'; t)$ est une algèbre d'Azumaya, et $k[u]$, $k[v]$ y sont des
sous-algèbres étales en position générale ; dans tous les cas c'est une algèbre
associative vérifiant (2), (3). La preuve (pages 88 à 90) réalise localement
$u$ et $v$ dans $M_2$ à partir de l'invariant $\rho_0$ : $\tau =
\operatorname{Tr}(uv)$ est alors une racine de $T^2 - bb'T + (b^2c' + b'^2c -
4cc' - \delta\delta'/\rho_0) = 0$, comme $t$ ; si $\tau \neq t$, $\tau = bb' - t$,
et l'on remplace $u$ par $u' = b - u$, pour lequel $\operatorname{Tr}(u'v) =
bb' - \tau = t$. « On gagne ! »

\emph{Remarque} (pages 90 et 91). Si $\delta\delta' \in k^{*}$ et si $M$ est
Azumaya, alors $Q \in k^{*}$, car $\rho_0 Q = \delta\delta'$ ; le discriminant de
$Q$ comme polynôme en $t$ est
\[
  (13)\qquad \Delta(Q) = (bb')^2 - 4(b^2c' + b'^2c - 4cc') = \delta\delta',
\]
de sorte que $2t - bb'$ est inversible et que le revêtement des racines de
$Q$ est étale. Le corollaire de la page 91, marqué « (?) » puis « Absurde »,
voudrait que $Q(u, v) = 0$ force $u$ et $v$ dans une même sous-algèbre de
Borel $k \oplus ku \oplus kv$ ; il est barré.

\pagerange{92}{99}
\subsection*{Le critère (pages 92 à 99)}

La matrice de la forme $\varphi$ sur la base $1, u, v, uv$ est
\[
  \begin{pmatrix}
    2 & b & b' & t \\
    b & b^2 - 2c & t & bt - b'c \\
    b' & t & b'^2 - 2c' & b't - bc' \\
    t & bt - b'c & b't - bc' & t^2 - 2cc'
  \end{pmatrix}
\]
(page 92)\footnote{La page porte $t^2 - cc'$ en bas à droite ; on attend
$\operatorname{Tr}((uv)^2) = t^2 - 2cc'$, comme le note la transcription.}. Sur
la base $1, u, v$ de $k1 + ku + kv$, le discriminant de $\varphi$ vaut $-2Q$,
et le discriminant réduit de la forme quadratique $\det$ vaut $-Q$ (page 94,
« Vérifier si c'est ça ? — Oui »). Si $Q$ n'est pas inversible, il existe
localement, sur une fibre, un $\Theta \neq 0$ dans $k1 + ku + kv$ orthogonal à
$1, u, v$ ; il est de trace nulle et isotrope, donc nilpotent, et les droites
isotropes de l'hyperplan de trace nulle correspondent aux Borels : $u$ et $v$
sont dans un même Borel. Le cas $\delta = \delta' = 0$ (sous-algèbres
unipotentes) se ramène à $u$, $v$ nilpotents, où $Q = t^2$ (pages 95 et
96)\footnote{Le crochet de la page 95 note qu'en caractéristique $2$ on ne
peut pas toujours ramener $b$ à $0$ quand $b^2 = 0$ : « ça coince ». La page
96 écrit que $\operatorname{Tr}(ux) = 0$ signifie que $x$ normalise le Borel
$B_u$, « i.e. $[x, u] = 0$ » ; la seconde formulation est fausse, la première
juste : pour $u$ nilpotent non nul, $u^{\perp} = B_u$ est la sous-algèbre de
Borel des éléments qui stabilisent la droite $\operatorname{Im}u$.}.

\textbf{Proposition} (page 97). $M(0, 0, 0, 0; t)$ est Azumaya si et seulement
si $t \in k^{*}$ ; on le réalise par $u = \bigl(\begin{smallmatrix} 0 & 1 \\ 0
& 0\end{smallmatrix}\bigr)$, $v = t\bigl(\begin{smallmatrix} 0 & 0 \\ 1 &
0\end{smallmatrix}\bigr)$.

\textbf{Théorème} (page 98). Soient $b, c, b', c', t \in k$ et $M = M(b, c; b',
c'; t)$. Pour que $M$ soit une algèbre d'Azumaya, il faut et il suffit que la
forme $\det$ sur $k \oplus ku \oplus kv$ soit non dégénérée, c'est-à-dire que
\[
  Q = t^2 - bb't + b^2c' + b'^2c - 4cc'
\]
soit inversible. Intrinsèquement : si $M$ est une algèbre localement libre de
rang $4$, munie de $\operatorname{Tr}$ et $\det$ avec les propriétés
habituelles, et $A_1$, $A_2$ deux sous-algèbres de rang $2$ transverses dont la
somme est un sous-module de rang $3$, $M$ est Azumaya si et seulement si $\det$
est non dégénérée sur $A_1 + A_2$.

La démonstration de la page 99 se place sur un corps algébriquement clos.
Nécessité : les cas « $A_1$, $A_2$ multiplicatifs » et « étales » sont acquis ;
il reste une position relative, $v$ unipotent et $u$ projecteur, $Q = t^2$, où
l'on choisit $u$ général. La suffisance n'est pas écrite (« Suffisance : … »).
Quand $2$ est inversible elle se déduit d'un calcul : le discriminant de la
forme $\varphi$ sur la base $1, u, v, uv$ de $M$ vaut
\[
  \det\bigl(\operatorname{Tr}(e_ie_j)\bigr) = -Q^2,
\]
de sorte que $Q$ inversible rend la forme trace non dégénérée, ce qui (pour une
algèbre de rang $4$ dont tout élément satisfait Hamilton--Cayley) fait de $M$
une algèbre d'Azumaya\footnote{Le calcul est de nous. En caractéristique $2$,
la forme trace est dégénérée et l'argument ne s'applique pas ; on ne sait pas
ce que la page comptait faire. La station X (pages 121 à 126) revient sur la
réalisabilité, et y constate que l'invariant $\rho$ était « une complication
inutile ».}.

\pagerange{101}{147}
\section*{X. Azumaya rang 4 (II) (pages 101 à 147)}

Datée « Avril 86 ». La suite des pages 104 à 108 se lit à rebours de la
numérotation : page 108 (proposition), 107 (démonstration), 106, 105, 104,
chaque feuillet continuant la phrase laissée en suspens au bas du
précédent\footnote{Ordre établi par la transcription.} ; de même la page 134 se
lit entre les pages 132 et 133, et la page 145 avant la page 144.

\pagerange{101}{103}
\subsection*{Les deux formes (pages 101 à 103)}

Soit $B \subset A$ un sous-fibré vectoriel ; si $B$ est de rang $3$, $D =
B^{\perp}$ est de rang $1$, et $B \leftrightarrow D$ est bijective, la forme
$\varphi$ étant non dégénérée. Pour $B \ni 1$, les orthogonaux de $B$ pour
$\varphi$ et pour $\psi$ coïncident : $\varphi(1, u) = \psi(1, u) =
\operatorname{Tr}u$, donc $B^{\perp} \subset A^{\natural}$, et pour $\xi \in
A^{\natural}$, $\psi(u, \xi) = -\varphi(u, \xi)$ (page 102). Les discriminants,
eux, coïncident au signe près :
\[
  \operatorname{disc}(\psi|_B) = (-1)^{\operatorname{rk}B - 1}\operatorname{disc}(\varphi|_B)
  \qquad (B \ni 1),
\]
égaux donc pour $B$ de rang $3$, le cas utilisé\footnote{La page 101 affirme
l'égalité « pour $B$ de rang quelconque », par réduction au cas universel
($2$ inversible, $B = k1 \oplus V$, $V \subset A^{\natural}$), en disant que
sur $V$ les deux formes « sont égales ». Elles sont opposées sur $V$, d'où le
signe. Pour $B = k1 + ku$ : $\operatorname{disc}\varphi = b^2 - 4c$,
$\operatorname{disc}\psi = 4c - b^2$. Pour le rang $3$, les deux valent $-2Q$.}.
La page 103 voit dans $B \mapsto \delta_B$ une section de
$\mathcal{O}(2)$ sur le plan projectif des hyperplans de $A/k$, venue de la forme
quadratique canonique $\det$ sur $A^{\natural} \simeq (A/k)^{\vee}$ : seuls
interviennent $\det$ et l'élément $1$, non la multiplication.

\pagerange{104}{109}
\subsection*{Sous-modules de rang $3$ (pages 104 à 109)}

Pour $B$ de rang $3$ contenant $1$, la multiplication induit une application
alternée $\mu : \bigwedge^2 B \to A/B$, $u \wedge v \mapsto uv \bmod B$ (alternée
par Hamilton--Cayley), nulle dès qu'un argument est dans $k$, donc une section
$\mu$ de $(A/B)^{\otimes 2} \simeq D^{\otimes -2}$ ; et $B$ est une sous-algèbre
si et seulement si $\mu = 0$. On pose $\delta(B) = \operatorname{disc}'(\det|_B)$,
le discriminant divisé, section de $(\det B)^{\otimes -2} \simeq
D^{\otimes -2}$, avec $\operatorname{disc}\varphi_B =
\operatorname{disc}\psi_B = 2\,\delta(B)$ ; c'est aussi $\det|_D$ (page 105).
Le point est que $\mu = \pm\delta(B)$, ce qu'il suffit de voir sur le cas
universel $\mathbb{P}^2_{\mathbf{Z}}$, où les deux sections ont les mêmes zéros
(page 104).

\textbf{Proposition} (page 108). Sur un corps (ou fibre à fibre), les
conditions suivantes sont équivalentes : (i) $D$ est isotrope ; (ii) $\delta(B)
= 0$ ; (iii) $B$ est une sous-algèbre ; (iii') $B$ est une sous-algèbre de
Borel. Elles entraînent (iv) $D \subset B$ et (ii') $\operatorname{disc}\varphi_B
= 0$, et leur sont équivalentes si $2$ est inversible\footnote{La page énonce
la proposition sur une base quelconque, puis doute (page 106) de (iii)
$\Rightarrow$ (i) hors d'un corps ou d'un anneau réduit, « le cas-test décisif
est celui où $k = k_0[t]/t^2$ … (l'histoire de cohomologie d'une algèbre de
Borel) ». La démonstration (page 107) : (i) $\Leftrightarrow$ (ii) parce que
$\det$ est non dégénérée ; isotrope entraîne $D \subset D^{\perp} = B$, la
réciproque demandant $2$ inversible (en caractéristique $2$, on a toujours $D
\subset B$) ; enfin les Borels correspondent aux droites isotropes de
$A^{\natural}$, $B_D = D^{\perp}$ étant le normalisateur de $D$ pour le
crochet.}.

\textbf{Corollaire} (page 109). Sont équivalentes : $D$ anisotrope sur chaque
fibre ; $\delta(B)$ inversible, c'est-à-dire $\det|_B$ non dégénérée ; $B$
n'est une sous-algèbre sur aucune fibre ; $B$ engendre $A$ comme algèbre ; et,
si $2$ est inversible, $D \not\subset B$ sur les fibres, $\varphi|_B$ non
dégénérée. Si $1, u, v$ est une base de $B$, elles équivalent aussi à : $1, u,
v, uv$ est une base de $A$, ou $uv \bmod B$ engendre $A/B$. C'est le théorème
de la page 98 sous forme intrinsèque.

\pagerange{109}{120}
\subsection*{Deux sous-algèbres quadratiques (pages 109 à 120)}

Soient $H$, $H' \subset A$ deux sous-fibrés de rang $2$ contenant $1$ (ce sont
des sous-algèbres, par Hamilton--Cayley), avec $H \cap H' = k$ sur les fibres ;
$B = H + H'$ est de rang $3$. Les orthogonaux $V$, $V'$ de $H$, $H'$ sont des
plans de $A^{\natural}$, et $D = V \cap V'$. Dans le plan projectif de
$A^{\natural}$ il y a la conique des droites isotropes ; $H$ est étale si et
seulement si la droite projective $\mathbb{P}(V)$ coupe la conique en deux points
distincts, et un $H$ quelconque correspond à un diviseur de degré $2$ de la
conique (page 110 et sa marge). On dit $H$, $H'$ \emph{en position générale}
si sur chaque fibre géométrique les deux paires de points sont disjointes :
il n'y a alors pas de droite isotrope orthogonale à $B$, et par la proposition
$\delta_B$ est inversible (page 111)\footnote{Les objets sont ceux du dossier
82 : la conique, les sous-algèbres spéciales comme diviseurs de degré $2$, les
Borels comme points de la conique (ses pages 34 et 46 à 55, 76 et 77).}.

Le rapport anharmonique $\lambda$ des deux paires (choisies localement) est
inversible ainsi que $1 - \lambda$, devient $\lambda^{-1}$ quand on échange les
points d'une paire, et ne dépend pas de l'ordre des deux paires ; $\rho =
\lambda + \lambda^{-1}$ est défini sur $k$, et $\rho_0 = \rho - 2$ y est
inversible (page 112). Sur un corps algébriquement clos, deux paires étales en
position générale sont conjuguées si et seulement si elles ont même $\rho_0$, et
tout $\rho_0$ inversible est réalisé (pages 112 et 113). Sur une base, le
stabilisateur de $(H, H')$ dans $\underline{\operatorname{Aut}}(A)$ est
$(\mathbf{Z}/2)_S$ quand $\rho_0 + 4$ est inversible, l'élément non trivial
induisant sur $H$ et $H'$ l'involution $x \mapsto x'$ ; les isomorphismes de
$(H, H')$ sur $(\overline{H}, \overline{H}')$ forment un torseur sous ce groupe,
qui peut être non trivial : il faut aussi que les revêtements étales $\varepsilon_H$,
$\varepsilon_{\overline{H}}$ soient isomorphes (page 114). Exemple (page 115) :
pour $H' = \overline{H}'$ donné par $\{0, \infty\}$ et $H$, $\overline{H}$ par
$Z^2 - bZ + c$, $Z^2 - \bar bZ + \bar c$, on a $\rho = b^2/c - 2$ ; avec $b =
\bar b = 0$, $\rho = \bar\rho = -2$ (le cas harmonique, exceptionnel) mais les
extensions ne sont isomorphes que si $\bar c/c$ est un carré. La question de la
réalisation Zariski-locale de $\rho_0$ reste ouverte (« je suppose que oui »).

\emph{Expression} (pages 116 à 120). En base $1, u, v$ :
\[
  \rho_0 = \frac{\delta_H\,\delta_{H'}}{Q},\qquad \delta_H = b^2 - 4c,\
  \delta_{H'} = b'^2 - 4c',\qquad Q = -\delta_B ,
\]
où $\delta_H = \operatorname{disc}(\varphi|_H)$ relativement à la base $1, u$,
et $\delta_B$ le discriminant divisé de $\det|_B$ relativement à $1, u, v$. Les
trois sont des sections de $(\det B)^{\otimes -2} \simeq (H/k)^{\otimes -2}
\otimes (H'/k)^{\otimes -2} \simeq D^{\otimes -2}$, isomorphisme qui ne dépend
plus d'aucun signe, et
\[
  \rho_0 = -\,\frac{\delta_H\,\delta_{H'}}{\delta_B}
\]
est intrinsèque\footnote{La page 118 encadre $\rho_0 = \delta_H\delta_{H'}/\delta_B$,
sans le signe ; la page 120 écrit $-\delta_B = Q$, et le calcul (discriminant
réduit de $\det$ sur $1, u, v$ égal à $-Q$, page 94) le confirme. La page 116
écrit $Q = \frac12\operatorname{disc}(\varphi|_B)$, avec devant un signe
noirci : c'est $-\frac12$.}. On peut le prendre pour définition dès que $\delta_B$
est inversible, et plus généralement comme section de $\mathbb{P}^1_k$ quand
$\delta_H\delta_{H'}$ et $\delta_B$ ne s'annulent pas ensemble : $\rho_0 =
\infty$ là où $H$, $H'$ sont dans un Borel, $\rho_0 = 0$ là où l'un n'est pas
étale (page 119). Sur un corps, le seul cas où $\rho_0$ n'est pas défini est
celui d'un Borel $B = B_D$ avec $H = k \oplus D$ et $H'$ un autre plan de $B$
contenant $k$ — situation standard, « couple de Killing » (un Borel et un tore
maximal dedans).

\pagerange{121}{126}
\subsection*{La réalisation (pages 121 à 126)}

La page 120 annonce, pour tout $(b, c, b', c', t)$ avec $Q$ inversible, une
réalisation dans une forme de $M_2$ par deux éléments en position relative
générale (la marge : « sans doute il suffit de l'inversibilité de $\delta$ »).
Sur un corps, l'argument de la page 90 revient (pages 121 et 122) : $Q(t) =
Q(\tau)$, donc $\tau \in \{t, bb' - t\}$, et l'on remplace au besoin $u$ par
$u'$. Le cas $\delta\delta' = 0$, « oublié », se ramène à trois cas types (page
123) : (1) $u$, $v$ idempotents, $Q = t(t - 1)$ ; (2) $u$ idempotent, $v^2 = 0$ ;
(3) $u^2 = v^2 = 0$ ; dans les cas 2 et 3, $Q = t^2$. Sur un corps
algébriquement clos, avec $A = \operatorname{End}(V)$, $u = x' \otimes x$, $v =
y' \otimes y$, on a $\operatorname{Tr}(uv) = \langle x, y'\rangle\langle y,
x'\rangle$ (page 124), et chaque cas se résout à la main : par exemple dans le
cas 1, $y_1y'_1 = t$, $y_2y'_2 = 1 - t$ (page 125). « Finalement, on n'a pas eu
besoin de faire intervenir $\rho$ dans la vérification, c'était une
complication inutile » (page 126).

\pagerange{126}{136}
\subsection*{La théorie intrinsèque : twists (pages 126 à 136)}

Soient $H$, $H'$ deux algèbres quadratiques sur $k$ (formes tordues de
$k[u]$, $k[v]$), avec leurs normes $\Delta_H$, traces $T_H$ et unités. La page
126 note que la forme $\varphi_H(u, v) = T_H(uv)$ se déduit de $\Delta_H$ et
$T_H$ ($\varphi_H + \psi_H = T_H \otimes T_H$), et la marge de la page 127 que
la structure d'algèbre de $H$ résulte déjà de la forme quadratique $\Delta_H$
et de l'élément $1_H$ avec $\Delta_H(1_H) = 1$. On forme
\[
  B = H \amalg_k H' = (H \oplus H')/k(1_H - 1_{H'}),\qquad P = H \otimes H',
\]
et l'on a une suite exacte $0 \to B \to P \to L \otimes L' \to 0$, avec $L =
H/k$, $L' = H'/k$ (page 134). Des plongements $H, H' \subset A$ dans une algèbre
d'Azumaya équivalent à une application linéaire $f : P \to A$, $u \otimes v
\mapsto uv$, dont les restrictions à $H$ et $H'$ sont des homomorphismes
(page 129 ; $f$ n'est pas supposée bijective, et $P$ n'est pas une algèbre).
Elle transporte sur $P$ une forme linéaire $L_P = \operatorname{Tr} \circ f$,
une forme quadratique $\Delta_P = \det \circ f$ avec $\Delta_P(u \otimes v) =
\Delta_H(u)\Delta_{H'}(v)$, et les formes $\psi_P$, $\varphi_P$.

\textbf{Proposition et théorème récapitulatif} (pages 130 à 136). Les ensembles
\begin{itemize}
  \item $\underline{Q}_P$ des formes quadratiques sur $P$ vérifiant
    $\Delta_P(u \otimes v) = \Delta_H(u)\Delta_{H'}(v)$,
  \item $\underline{Q}_B$ des formes quadratiques sur $B$ prolongeant
    $\Delta_H$, $\Delta_{H'}$,
  \item $\underline{\Psi}_B$, $\underline{\Phi}_B$ des formes bilinéaires
    symétriques sur $B$ prolongeant $\psi_H, \psi_{H'}$, resp. $\varphi_H,
    \varphi_{H'}$,
  \item $\underline{L}_P$ des formes linéaires sur $P$ prolongeant $T_B$,
    c'est-à-dire des formes bilinéaires $T$ sur $H \times H'$ avec $T(u, 1) =
    T_H(u)$, $T(1, v) = T_{H'}(v)$,
  \item $AZ_P$ des structures d'algèbre de Hamilton--Cayley sur $P$ induisant
    les structures données sur $H \otimes 1$ et $1 \otimes H'$,
\end{itemize}
sont des torseurs sous $(L \otimes L')^{\vee}$, et les flèches naturelles
\[
  AZ_P \to \underline{Q}_P \to \underline{Q}_B \to \underline{\Psi}_B \to
  \underline{\Phi}_B \to \underline{L}_P
\]
(restriction, forme polaire, $\varphi = T_B \otimes T_B - \psi$, $\varphi
\mapsto (x \otimes y \mapsto \varphi(x, y))$) sont des isomorphismes de
torseurs, le diagramme qu'elles forment avec $AZ_P \to \underline{L}_P$,
$\alpha \mapsto \operatorname{Tr}_\alpha$, étant commutatif\footnote{La page 131
justifie la structure de torseur : les formes bilinéaires symétriques sur $B$
nulles sur $H$ et $H'$ sont les $((x, y), (\xi, \eta)) \mapsto \lambda(x \otimes
\eta) + \lambda(\xi \otimes y)$, $\lambda \in (L \otimes L')^{\vee}$. Le lemme de la
page 134 : $P \simeq \underline{\operatorname{Hom}}(\check{H}, H') \to
\underline{\operatorname{Hom}}(\det\check{H}, \det H')$ est une fonction
quadratique canonique, universelle parmi celles qui s'annulent sur les tenseurs
décomposables (page 133). Les marges des pages 127 et 136 notent que
$\underline{L}_P$ ne dépend que des $T$ et des unités, alors que la flèche
$\delta$ ci-dessous fait appel aux normes.}. Un élément $T_P \in \underline{L}_P$
est appelé un \emph{twist} des algèbres quadratiques $H$, $H'$ ; c'est l'objet
le plus simple, et la structure d'algèbre qu'il définit sur $P$ est « le plus
fondamental ». En outre : la structure d'algèbre associée à un twist est
d'Azumaya si et seulement si $\Delta_B$ est non dégénérée, c'est-à-dire
$\operatorname{disc}'(\Delta_B) \in (L \otimes L')^{\otimes -2}$ inversible ; et
$\delta_{P,B} : \underline{L}_P \to (L \otimes L')^{\otimes -2}$ est une
application polynomiale de degré $2$ dont la partie quadratique est
l'élévation au carré, donc un morphisme fini surjectif de degré $2$.

En coordonnées, $\underline{L}_P$ est l'axe des $t$, l'algèbre associée à $t$
est $M(b, c; b', c'; t)$, et $\delta_{P,B}$ est $t \mapsto Q(t)$ au signe près :
c'est la description intrinsèque de la station IX.

\pagerange{137}{147}
\subsection*{Le produit contracté $X \wedge X'$ (pages 137 à 147)}

Soient $X = \operatorname{Spec}H$, $X' = \operatorname{Spec}H'$ les revêtements
quadratiques de $S = \operatorname{Spec}k$. L'image réciproque
$\widetilde{S} = \delta_{P,B}^{-1}(0)$ est un revêtement quadratique de $S$, et
il y a un morphisme canonique $X \times_S X' \to \widetilde{S}$ compatible aux
involutions (page 137) ; si $X$ et $X'$ sont étales, $\widetilde{S}$ est le
produit contracté $X \wedge X'$ des $(\mathbf{Z}/2)_S$-torseurs. Le morphisme
s'écrit : à deux caractères $\varphi : H \to k$, $\varphi' : H' \to k$ on
associe le twist
\[
  t_{\varphi, \varphi'}(u, v) = \operatorname{Tr}(v)\varphi(u) +
  \operatorname{Tr}(u)\varphi'(v) - 2\varphi(u)\varphi'(v),
\]
qui est bien dans $\underline{L}_P$ ($t(u, 1) = \operatorname{Tr}u$, $t(1, v)
= \operatorname{Tr}v$), et qui annule $\delta$ : en coordonnées, avec $z =
\varphi(u)$, $z' = \varphi'(v)$ racines de $z^2 - bz + c$, $z'^2 - b'z' + c'$,
$t = b'z + bz' - 2zz'$ est racine de $Q$, l'autre racine étant $bb' - t$ (pages
138 et 139)\footnote{Vérifié. La page 138 écrit d'abord le terme constant de $Q$
comme $b^2b' + b'^2b - 4cc'$, et $(2z' - b)^2$ pour $(2z' - b')^2$ ; elle trouve
la formule par l'heuristique $2t - bb' = \pm(2z - b)(2z' - b')$ en
caractéristique différente de $2$, et ajoute que la vérification de $tt' =
b^2c' + b'^2c - 4cc'$ est « immédiate ». Géométriquement, $\delta(t_{\varphi,
\varphi'}) = 0$ dit que pour ce twist, $H$ et $H'$ ont un point fixe commun.}.
L'involution canonique $T \mapsto T'$, $T'(u, v) = T(u', v) = T(u, v')$, de
$\underline{L}_P$ (page 140 : $\operatorname{Tr}(u'v) = \operatorname{Tr}(uv')$
dans l'algèbre associée) s'écrit $t \mapsto bb' - t$ et induit l'involution
canonique de $\widetilde{S}$ (page 141).

\textbf{Proposition} (pages 141 et 142). Si $2$, $\delta$, $\delta'$ engendrent
l'idéal unité (en tout point où $2$ est nul, $X$ ou $X'$ est étale), $X \times_S
X' \to X \wedge_S X'$ est fini localement libre de rang $2$ et identifie $X
\wedge_S X'$ au quotient de $X \times_S X'$ par $\sigma_X \times \sigma_{X'}$ ;
il est étale si et seulement si en tout point $X$ ou $X'$ l'est, et $X \wedge
X'$ est étale si et seulement si $X$ et $X'$ le sont. L'algèbre de $X \wedge
X'$ est la sous-algèbre $H \wedge H' = (H \otimes H')^{\sigma_H \otimes
\sigma_{H'}}$ de l'algèbre commutative $H \otimes H'$, de base $1$ et
\[
  u * v = \operatorname{Tr}(v)\,u + \operatorname{Tr}(u)\,v - 2\,u \otimes v,
\]
avec, dans l'algèbre quadratique $H \wedge H'$,
\[
  \operatorname{Tr}(u * v) = bb',\qquad \det(u * v) = b^2c' + b'^2c - 4cc',
  \qquad \operatorname{disc}(u * v) = \delta\delta' ,
\]
$t_{\varphi, \varphi'} = (\varphi \otimes \varphi')(u * v)$, et $(u * v)' = u' *
v = u * v'$\footnote{La page 143 écrit $\det(u * v) = b^2c + b'^2c' - 4cc'$, les
accents mal visibles ; la valeur juste est le terme constant de $Q$, puisque
$u * v$ annule $Q$. Le produit $H \wedge H'$ est le produit des algèbres
quadratiques que l'on trouve chez Knus (\emph{Quadratic and Hermitian Forms
over Rings}), et c'est l'opération $A_1 * A_2$ du dossier 84 (ses pages 62 à
106). Noms et rapprochement de nous.}. On récupère $\underline{L}_P$ comme
l'espace affine enveloppant de $\widetilde{S}$, et $\delta$ comme son
discriminant (page 143).

Le cas unipotent (pages 144 à 146) : si $H = k \oplus L$, $H' = k \oplus L'$
avec $L^2 = L'^2 = 0$, $H \wedge H' \simeq k \oplus (L \otimes L')$ et $X \times
X' \to X \wedge X'$ est $(\xi, \eta) \mapsto \xi \otimes \eta$ ; on a $u * v =
-2uv$, nul en caractéristique $2$, où $H \wedge H'$ ne se plonge plus dans $H
\otimes H'$. Ainsi $X_L \wedge X_{L'} \simeq X_{L \otimes L'}$ pour $X_L =
\operatorname{Spec}(k \oplus L)$, et $X \wedge X$ n'est pas trivial : le
produit contracté des revêtements quadratiques se comporte autrement que celui
des torseurs, sauf dans le cas étale. Quand $2$, $b$, $b'$ engendrent l'idéal
unité, $P = H \otimes H'$ contient trois sous-algèbres quadratiques $H$, $H'$,
$H'' = H \wedge H'$, fixées par trois involutions commutantes, deux à deux
d'intersection $k$ ; $P$ est engendré par $H$ et $H'$, et aussi par $H$ et
$H''$ et par $H'$ et $H''$ si et seulement si $H$ et $H'$ sont étales — contre-exemple
$u^2 = u$, $v^2 = 0$, $w = u * v = v - 2uv$, où $v^2 = w^2 = vw = 0$ (page 145).
La page 147 voudrait un produit contracté multiple $\bigwedge_{i \in I} X_i$
décrit par des formes multilinéaires $f$ avec $f(1, u_2, u_3) =
\operatorname{Tr}u_2\operatorname{Tr}u_3$, etc., formant un torseur sous $(L_1
\otimes L_2 \otimes L_3)^{-1}$, « satisfaisant à des conditions quadratiques
convenables — mais lesquelles ?? ». C'est l'objet de la suite.

\pagerange{149}{180}
\section*{XI. Azumaya rang 4 (III) : espaces affines à involution (pages 149 à 180)}

Datée « Avril 86 », ouverte par « Digression sur les produits tensoriels
d'espaces affines à involution ».

\pagerange{149}{160}
\subsection*{Le produit $*$ (pages 149 à 160)}

Sur un module $H$ muni d'une unité ($k \to H$ injective), une application
linéaire $x \mapsto x'$ qui fixe $1$ et induit $-1$ sur $H/k$ équivaut à une
forme linéaire $T$ avec $T(1) = 2$, par $x' = T(x)1 - x$. Dualement, un espace
affine $A = \varepsilon^{-1}(1)$ d'enveloppe $\varepsilon : E \to k$ (module des
translations $A^0 = \ker\varepsilon$), muni d'une involution induisant
$-\mathrm{id}$ sur $A^0$, équivaut à un $t \in E$ avec $\varepsilon(t) = 2$,
l'involution étant $x \mapsto \varepsilon(x)t - x$. Si $2$ est inversible elle a
un unique point fixe et tout est trivial ; si $2 = 0$, $t \in A^0$ et
l'involution est la translation par $t$ (pages 149 et 150).

Une application $f : A_1 \times A_2 \to B$ est \emph{$(\sigma_1,
\sigma_2)$-bilinéaire} s'il existe $f^0 : A_1^0 \times A_2^0 \to B^0$ bilinéaire
avec
\[
  f(x_1 + \alpha_1, x_2) = f(x_1, x_2) + f^0(\alpha_1, \sigma_2x_2 - x_2),\qquad
  f(x_1, x_2 + \alpha_2) = f(x_1, x_2) + f^0(\sigma_1x_1 - x_1, \alpha_2),
\]
d'où $f(x_1 + \alpha_1, x_2 + \alpha_2) = f(x_1, x_2) + f^0(\alpha_1, \tau_2) +
f^0(\tau_1, \alpha_2) - 2f^0(\alpha_1, \alpha_2)$ en coordonnées ($\tau_i =
\sigma_ix_i - x_i$)\footnote{La page 151 écrit $f^0(\sigma_1x_1 - x_1, \alpha_1)$
au troisième terme ; c'est $\alpha_2$. Comme $f$ ne détermine pas toujours
$f^0$ quand $2$ n'est pas inversible (si $2 = 0$ et $\tau_i = 0$, $f$ est
constante), la page met $f^0$ dans la donnée (page 152).}. Il y a une solution
universelle $A_1 * A_2$, espace affine sous $A_1^0 \otimes A_2^0$, réalisée
dans $E_1 \otimes E_2$ par
\[
  x_1 * x_2 = t_1 \otimes x_2 + x_1 \otimes t_2 - 2\,x_1 \otimes x_2 = \sigma_1(x_1)
  \otimes x_2 + x_1 \otimes \sigma_2(x_2),
\]
avec $\sigma_1(x_1) * x_2 = x_1 * \sigma_2(x_2) = t_1 \otimes t_2 - x_1 * x_2$ :
les deux involutions induites coïncident, et $A_1 * A_2$ est encore un espace
affine à involution, solution du problème universel dans cette catégorie (page
153). L'opération est associative et commutative (page 154 et sa marge) ; son
unité est la droite $k$ avec l'involution $x \mapsto 1 - x$ (et non $x \mapsto
-x$, triviale en caractéristique $2$), via $x \mapsto e_0 * x$, $e_0 = (1,
0)$. Pour $n$ facteurs, le produit est défini par le problème universel analogue
et réalisé dans $\bigotimes E_i$ par
\[
  x_1 * \cdots * x_n = \sum_{\substack{J \subset \{1, \ldots, n\} \\ \operatorname{card}J \text{ impair}}}
  \ \bigotimes_{j \in J} x_j \otimes \bigotimes_{j \notin J} \sigma_j(x_j)
  \qquad (2^{n-1} \text{ termes}),
\]
l'involution du produit donnant la même somme sur les $J$ de cardinal pair,
et $x_1 * \cdots * x_n + \sigma(x_1 * \cdots * x_n) = t_1\cdots
t_n\,\varepsilon_1(x_1)\cdots\varepsilon_n(x_n)$ (page 156)\footnote{La page
155 donne aussi la forme développée $\sum_i (-2)^{i-1}(\sum t\cdots t\,x\cdots
x)$, avec un signe et un coefficient final surchargés, et la page 156 l'image
dans $\bigotimes E_i$. L'enveloppe de $\mathop{*}A_i$, quand $2$ est inversible,
a pour augmentation $\varepsilon(t_1\cdots t_n) = 2$, et non $2^n$ (page 157).}.

Si $2$ est inversible, $\bigotimes E_i$ se décompose selon les caractères de
$(\mathbf{Z}/2)^n$, l'enveloppe du produit est le sous-espace propre du
caractère $(1, \ldots, 1)$, et l'opération n'est que le produit tensoriel fait
séparément sur $E^-$ et $E^+$ — « aucun mystère ! » ; elle n'est « délicate et
utile » que si $2$ n'est pas inversible (pages 157 à 160). Pour des modules
involutifs quelconques $(E_i, \sigma_i)$ avec $E_i^0 \subset E_i$, le produit
universel $\mathop{*}(E_i, \sigma_i)$ est une extension de $\bigotimes(E_i/E_i^0)$
par $\bigotimes E_i^0$, et l'application $\tau$ vers $\bigotimes E_i$ n'est en
général ni injective ni surjective (page 159).

\pagerange{160}{170}
\subsection*{Revêtements quadratiques et le torseur $Q(A)$ (pages 160 à 170)}

Une droite affine $A$ sous un module inversible $L$ et un diviseur
quadratique de $\check{\mathbb{V}}(A)$ transverse à l'infini équivalent à une
forme quadratique $q : E \to L^{\otimes 2}$ dont la partie homogène est $u
\mapsto u^{\otimes 2}$ ; c'est la donnée d'un revêtement quadratique de $S$
(pages 160 et 161). Elle définit une involution de $E$ compatible à
l'augmentation, qui laisse le diviseur stable et y induit l'involution canonique
(page 162)\footnote{La page 162 met en garde : en caractéristique $2$,
l'invariance par une involution donnée ne suffit pas à caractériser les
diviseurs. Les conventions de signe (« Attention aux conv. de signe ! ») ne sont
pas fixées sur les pages.}. Les formes quadratiques compatibles avec une
involution donnée forment un torseur sous $L^{\otimes 2}$, identifié à $A * A$
via $\xi + q_\xi(x) = x * x$ (pages 163 et 164). La page 165 construit le
torseur universel $Q(A, \sigma)$ sous $L^{\otimes 2}$ et la fonction
sous-quadratique $A \to Q(A)$, revêtement de degré $2$, dont les sections
correspondent aux revêtements quadratiques d'involution donnée.

Pour $A = A_1 * A_2$, en coordonnées (origines $o_i$, $\sigma_i(x) = b_i -
x$), les formes compatibles s'écrivent $q_{c_i}(x) = x^2 - b_ix + c_i$, et
\[
  q_{c_1} \otimes q_{c_2} = q_c,\qquad c = b_2^2c_1 + b_1^2c_2 - 4c_1c_2,
\]
le terme constant de $Q$ ; d'où $(q_1 + \alpha) \otimes q_2 = q_1 \otimes q_2 +
\alpha\,\delta(q_2)$ et $\delta(q_1 \otimes q_2) = \delta(q_1)\delta(q_2)$, et un
diagramme commutatif reliant $A_1 \times A_2 \to A$, $Q_1 \times Q_2 \to Q$ et
les discriminants (pages 168 à 170). La page voudrait écrire $Q(A_1 * A_2) \simeq
Q(A_1) * Q(A_2)$, mais les $Q(A_i)$ ne sont pas à involution : « il faudrait
reprendre la formalisation de façon plus générale ».

\pagerange{171}{180}
\subsection*{Extensions $\lambda$-scindées (pages 171 à 180)}

Soit $0 \to E^{\circ} \to E \xrightarrow{\varepsilon} F \to 0$ une extension. Une
involution de $E$ induisant $-1$ sur $E^{\circ}$ et $+1$ sur $F$ s'écrit $\sigma
= \pi\varepsilon - \mathrm{id}$ avec $\varepsilon\pi = 2\,\mathrm{id}_F$ : $\pi$ est
une « $2$-section ». Plus généralement, pour $\lambda \in k$, une
\emph{$\lambda$-section} $\pi : F \to E$, $\varepsilon\pi = \lambda\,\mathrm{id}_F$,
est un scindage de l'extension image réciproque par $\lambda$ ; elle équivaut à
une $\lambda$-rétraction
\[
  u^{\pi} : E \to E^{\circ},\qquad u^{\pi}(x) = \lambda x - \pi\varepsilon(x),
  \qquad u^{\pi}|_{E^{\circ}} = \lambda\,\mathrm{id},
\]
et c'est $u$ qui sert : pour $\lambda = 2$ et $F = k$, $u(x) = 2x - t = -(\sigma(x)
- x)$, l'expression qui jouait le rôle crucial dans la construction de $*$
(pages 171 à 174). Un \emph{$\lambda$-accouplement} est un couple $(f, f^{\circ})$
d'applications multilinéaires avec
\[
  f(\ldots, x_i + \alpha_i, \ldots) = f(\ldots, x_i, \ldots) + f^{\circ}(u_1(x_1),
  \ldots, \alpha_i, \ldots, u_n(x_n)),
\]
d'où $f(\alpha_1, \ldots, \alpha_n) = \lambda^{n-1}f^{\circ}(\alpha_1, \ldots,
\alpha_n)$ ; la solution universelle $\mathop{*}_{i, \lambda} E_i$ est un quotient de
$\bigotimes E_i \oplus \bigotimes E_i^{\circ}$, avec une suite exacte
$\bigotimes E_i^{\circ} \to \mathop{*}_{i,\lambda} E_i \to \bigotimes F_i \to 0$,
injective à gauche si les $F_i$ sont plats (pages 175 à 177). La
$\lambda$-rétraction canonique du produit vient du $\lambda$-accouplement $\bigotimes
u_i$ (page 177). On obtient une opération interne dans les extensions
$\lambda$-scindées (à quotients plats), associative et commutative « à
vérifier » ; la question de l'unité (pages 179 et 180) : pour l'extension
scindée $E_\beta$ de $k$ par $k$ de $\lambda$-trivialisation $\beta$, $E_\beta
*_\lambda E_2 \simeq {}^{\beta}E_2$, qui est une équivalence si et seulement si
$\beta$ est inversible, et l'identité pour $\beta = 1$. La page 180 énonce trois
conditions équivalentes pour qu'un $M$ $\lambda$-trivialisé soit inversible pour
$*_\lambda$ ($M^{\circ}$ et $M/M^{\circ}$ inversibles, $u : M \to M^{\circ}$
surjectif, …) et s'arrête\footnote{Ce sont les $\lambda$-rétractions et le produit
des morphismes du dossier 84, troisième partie (ses pages 123 à 191),
notamment ses pages 133 à 139.}.

\pagerange{182}{200}
\section*{XII. (I bis) : la trace d'un produit de trois matrices (pages 182 à 200)}

Titre cerclé « (I bis) », marge « Mai 1986 — Juillet 1986 ». Les pages 194,
196 et 198 à 200 sont écrites au verso de pages dactylographiées.

\pagerange{182}{191}
\subsection*{Trois variables (pages 182 à 191)}

On veut expliciter la multiplication de $A$ à partir de $\det$, de $1$ et de
l'orientation $e \in \bigwedge^4 A$ associée à $\det$ ($e \otimes e =
\operatorname{disc}(\det)$, une « structure spéciale orthogonale »). Comme
$\psi$ est connue et non dégénérée, $xy$ est déterminé par les
$\operatorname{Tr}(xyz)$ ; mais $\det$, $1$ et $e$ ne fixent la multiplication
qu'au passage près à l'opposée, c'est-à-dire à l'involution $x \mapsto x'$ près,
et l'on attend donc une ambiguïté. Avec les notations $b, b', b''$, $c, c',
c''$, $t = \operatorname{Tr}yz$, $t' = \operatorname{Tr}zx$, $t'' =
\operatorname{Tr}xy$ et $\tau = \operatorname{Tr}(xyz)$, $\tau' =
\operatorname{Tr}(zyx)$ (chacune invariante par permutation circulaire), (6')
donne
\[
  \tau + \tau' = \Sigma = bt + b't' + b''t'' - bb'b'' .
\]
La différence $\tau - \tau' = \operatorname{Tr}([x, y]z)$ est trilinéaire
alternée en $x, y, z$ et invariante par $\underline{\operatorname{Aut}}(A)$ ; elle
est donc de la forme $\varepsilon\,(1 \wedge x \wedge y \wedge z) \otimes e^{-1}$
(page 185), et le triplet $x_0 = E_{21}$, $y_0 = E_{12}$, $z_0 = [x_0, y_0]$
fixe $\varepsilon$ :
\[
  \tau - \tau' = \varepsilon_0\,\Delta,\qquad \Delta = (1 \wedge x \wedge y \wedge
  z) \otimes e^{-1},\qquad \varepsilon_0 = \pm 1,
\]
le signe dépendant de la convention d'orientation\footnote{La page 184 écrit
d'abord $2(\tau - \tau') = \varepsilon_0\Delta$, avec un coefficient surchargé ;
la page 188 a $\tau - \tau' = \varepsilon_0\Delta$, qui est juste (vérifié). La
page 185 écrit $\tau' - \tau$ pour $\operatorname{Tr}xyz -
\operatorname{Tr}yxz$, sans effet sur la suite.}. Une première tentative
d'écrire $\Delta^2$ comme polynôme en les $b$ et les $t$ seuls échoue sur
l'exemple ($4 = \alpha\cdot 0$, « ce qui est fâcheux ! ») ; « une fois qu'on y
pense », $\Delta^2$ est le déterminant de Gram de $\psi$ sur $1, x, y, z$ :
\[
  \Delta^2 = \det\begin{pmatrix}
    2 & b & b' & b'' \\
    b & 2c & bb' - t'' & bb'' - t' \\
    b' & bb' - t'' & 2c' & b'b'' - t \\
    b'' & bb'' - t' & b'b'' - t & 2c''
  \end{pmatrix}
\]
(page 187). D'où $4\tau\tau' = \Sigma^2 - \Delta^2$ et
\[
  \tau\tau' = \Pi = tt't'' + (ct^2 + c't'^2 + c''t''^2) + (b^2c'c'' + b'^2c''c +
  b''^2cc') - (b'b''ct + b''bc't' + bb'c''t'') - 4cc'c''
\]
(page 188)\footnote{La page 188 omet le terme $-4cc'c''$ ; sans lui l'identité
$\tau\tau' = \Pi$ est fausse dès que $cc'c'' \neq 0$. Nous avons déterminé
$\Pi$ en résolvant l'identité $\operatorname{Tr}(xyz)\operatorname{Tr}(zyx) =
\Pi$ sur les matrices génériques ; le polynôme obtenu est unique. Le $\Pi$ de
la page coïncide avec celui-ci pour des triplets dont l'un des éléments est
singulier. C'est l'identité de traces classique de Fricke ; nom de nous.}.

\emph{À vérifier} (pages 189 à 191) : l'anneau $B$ des fonctions polynomiales
$G$-invariantes de trois variables dans $A$ ($G = \underline{\operatorname{Aut}}(A)
\simeq A^{*}/\mathbb{G}_m$) serait libre de rang $2$ sur l'anneau de polynômes
$k[\underline{b}, \underline{b}', \underline{b}'', \underline{c}, \underline{c}',
\underline{c}'', \underline{t}, \underline{t}', \underline{t}'']$, engendré par
$\tau$ soumis à $\tau^2 - \Sigma\tau + \Pi = 0$ ; son involution canonique est
$\tau \mapsto \tau' = \Sigma - \tau$, et $\mathfrak{S}_3$ y opère, les
permutations circulaires fixant $\tau$, les transpositions échangeant $\tau$ et
$\tau'$\footnote{La page 189 écrit $\underline{c}'' = \det y$ pour $\det z$.
Sur un corps de caractéristique nulle, c'est la description classique de
l'anneau des invariants de trois matrices $2 \times 2$ par conjugaison (Sibirskii,
Procesi, Formanek) ; nous n'avons pas vérifié l'énoncé sur une base
quelconque, que la page présente comme à vérifier.}. Si $2$ est inversible on
peut prendre $\Delta = \tau - \tau'$ pour générateur, avec $\Delta^2 = \Sigma^2 -
4\Pi$ ; en caractéristique $2$, $\Delta \equiv \tau + \tau' = \Sigma$ est dans
l'anneau de base et n'engendre rien, donc pas non plus si $k/2k \neq 0$ (page
191).

\pagerange{191}{200}
\subsection*{Quatre variables, et le cube (pages 191 à 200)}

Pour quatre variables $x_1, \ldots, x_4$, on a les $b_i$, $c_i$, $t_{ij}$, soit
$4 + 4 + 6 = 14$ fonctions sur un quotient de dimension $16 - 3 = 13$ : il doit
y avoir une relation. Les $t_{ijk} = \operatorname{Tr}(x_ix_jx_k)$ sont $8$
distincts (invariance circulaire), groupés par paires $t_{ijk}$, $t_{jik}$ de
somme $\Sigma_{\{i,j,k\}}$ et de produit $\Pi_{\{i,j,k\}}$ (page 192). La page
193 présume l'anneau des invariants fini et libre sur l'anneau $A$ engendré par
les $b$, $c$, $t_{..}$, et voit la clef de la relation cherchée dans
$\operatorname{Tr}(x_1x_2x_3x_4)$. Les pages 196 et 198 notent que
$\operatorname{Tr}(u^n) = P_n(\operatorname{Tr}u, \det u)$, que $\tau =
\operatorname{Tr}(x_1x_2x_3x_4)$ satisfait, vu comme $\operatorname{Tr}(x_1x_2(x_3x_4))$,
une équation du second degré à coefficients dans les invariants de
$x_1, x_2, x_3x_4$, donc est entier de degré au plus $8$ sur $A$ ; la page 194
dessine le treillis des sous-anneaux fixes sous $\mathfrak{S}_3 \times
\mathfrak{S}_2$, la page 199 celui des degrés $2$ et $3$\footnote{Lecture mince :
ces pages sont des brouillons dispersés, en partie barrés, dont l'ordre de
lecture n'est pas sûr.}.

C'est la page 198 qui dessine le cube. La rotation du cube par $\mathfrak{S}_4$,
qui permute ses quatre grandes diagonales, organise les quantités : la
diagonale $\delta(s)$ d'un sommet $s$ porte $b_s$, $c_s$ ; une arête, dont les
extrémités sont sur deux diagonales, porte $t_a$ ; un sommet orienté porte un
$t_{ijk}$, une face orientée un $\operatorname{Tr}(x_ix_jx_kx_l)$ (pages 198 et
204). La page 200 en tire des relations : pour deux faces $f$, $f'$
adjacentes le long d'une arête $a$, $\tau_f + \tau_{f'} = \Sigma_a$, d'où des
relations linéaires $\Sigma_a + \Sigma_{\bar a'} = \Sigma_{a'} + \Sigma_{\bar
a}$, « douze relations » que $\mathfrak{S}_4$ permute\footnote{La page 200
affirme aussi que $\Sigma_a$ ne dépend que de la direction de $a$ ; la suite
(page 202) traite douze quantités $\Sigma_a$ distinctes et ne retient pas
l'affirmation.}.

\pagerange{202}{220}
\section*{XIII. Relations sur le cube (pages 202 à 220)}

Ces pages alternent avec des pages dactylographiées de \emph{Récoltes et
Semailles} ; elles sont des calculs presque sans prose, et la lecture qui suit
en donne la structure plus que le détail.

\pagerange{202}{206}
\subsection*{Variables et relations (pages 202 à 206)}

Les sommets $s$ du cube, ses arêtes $a$, ses faces $f$ portent (page 204)
\[
  b_s,\ c_s\ (s \in S),\qquad t_a\ (a \in A),\qquad \mathbb{T}_s =
  \operatorname{Tr}X_iX_jX_k,\qquad \mathbb{T}_f = \operatorname{Tr}X_iX_jX_kX_l,
\]
avec $b_s = b_{\bar s}$, $c_s = c_{\bar s}$, $t_a = t_{\bar a}$ pour les
éléments antipodaux ($4 + 4 + 6$ quantités), $\mathbb{T}_{\bar s}$ étant la
trace du produit en ordre inverse ($8$ quantités) et $f \mapsto (i, j, k, l)$
la correspondance entre faces et ordres circulaires de $I$ ($6$ quantités). Les
relations sont
\[
  (R_s)\quad \mathbb{T}_s + \mathbb{T}_{\bar s} = \sigma_s,\quad
  \mathbb{T}_s\mathbb{T}_{\bar s} = \pi_s ;\qquad
  (R_a)\quad \mathbb{T}_f + \mathbb{T}_{f'} = \sigma_a,\quad
  \mathbb{T}_f\mathbb{T}_{f'} = \pi_a ,
\]
où $\sigma_s$, $\pi_s$ sont les $\Sigma$, $\Pi$ de la station précédente
écrits sur les voisins de $s$\footnote{Le $\pi_s$ de la page 204 reprend le
$\Pi$ de la page 188 et hérite de son terme manquant $-4c\,c'c''$.}, et, pour
$f = (i, j, k, l)$, $f' = (j, i, k, l)$ adjacentes le long d'une arête,
\[
  \sigma_a = \operatorname{Tr}(X_kX_l)\operatorname{Tr}(X_iX_j) + \operatorname{Tr}X_i
  \operatorname{Tr}(X_jX_kX_l) + \operatorname{Tr}X_j\operatorname{Tr}(X_kX_lX_i) -
  \operatorname{Tr}X_i\operatorname{Tr}X_j\operatorname{Tr}(X_kX_l)
\]
(c'est (6') appliqué à $X_iX_j + X_jX_i$). Bilan de la page 206 : $28$
variables, $32$ relations ; en marge de la page 202, « Trop lourd ».

\pagerange{202}{220}
\subsection*{Les relations linéaires entre les $\Sigma_a$ (pages 202, 208 à 220)}

Puisque $\Sigma_a = \mathbb{T}_f + \mathbb{T}_{f'}$, les douze $\Sigma_a$ sont les
sommes, le long des douze arêtes du graphe d'adjacence des faces (un
octaèdre), de six quantités portées par les faces ; elles vérifient donc six
relations linéaires indépendantes\footnote{Explicitation de nous : l'octaèdre
ayant des cycles impairs, l'application des faces vers les arêtes est
injective, d'image de rang $6$.}. La page 202 en écrit douze (deux groupes de
six), observe que chacune des six dernières résulte des cinq autres, et note
qu'il existe $\sigma$ avec $\sigma_a + \sigma_b + \sigma_c = \sigma$ pour tout
« demi-hexagone » de trois arêtes deux à deux non coplanaires. Les pages 208 à
216 paramètrent les solutions : différences égales le long des faces
($\gamma' - c = a' - \beta' = \ldots = \delta'$), trois paramètres $\delta,
\delta', \delta''$ avec
\[
  \delta' + \delta'' = a' - a,\quad \delta'' + \delta = b' - b,\quad \delta +
  \delta' = c' - c,
\]
et les douze valeurs exprimées comme demi-sommes de $a \pm a'$, $b \pm b'$, $c
\pm c'$ (page 216). Les sommes $\varphi$ sur les arêtes des faces donnent une
application linéaire de $(a, b, c, a', b', c')$, dont la page 218 calcule le
déterminant pour conclure qu'il est non nul et premier à $3$. Ce déterminant
vaut\footnote{La page calcule $\det M = 27$, puis traite les blocs comme des
scalaires et trouve $\det\bigl(\begin{smallmatrix} 27 & 1 \\ 1 &
27\end{smallmatrix}\bigr) = 728 = 2^3\cdot 7\cdot 13$, « Ok ». Les blocs commutent,
le déterminant est $\det(M^2 - I)$ ; $M$ a pour valeurs propres $3, -3, -3$, donc
$M^2 - I = 8I$. La conclusion (non nul, premier à $3$) subsiste ; l'application
est inversible sur $\mathbf{Z}[\frac12]$, ce qui s'accorde avec les $\frac12$ des
pages 214 et 216. Le déterminant $6 \times 6$ de la page 210, qu'elle commence à
développer sans conclure, vaut $1024 = 2^{10}$ (notre calcul).}
\[
  \det\begin{pmatrix} M & I_3 \\ I_3 & M\end{pmatrix} = \det(M^2 - I_3) = 2^9 = 512,
  \qquad M = \begin{pmatrix} -1 & 2 & 2 \\ 2 & -1 & 2 \\ 2 & 2 & -1\end{pmatrix}.
\]
La page 220 pose
$\lambda = a + a'$, $\mu = b + b'$, $\nu = c + c'$, un paramètre par direction
d'arêtes, et s'arrête.

\pagerange{221}{221}
\section*{XIV. Le feuillet « L 19 » (page 221)}

Le dernier feuillet n'est pas de mathématiques : c'est une page
dactylographiée de sa main, paginée « L 19 », corrigée à l'encre bleue,
écrite autour de \emph{Récoltes et Semailles} et de « l'écriture des Quatre
Opérations ». Il y constate, depuis son départ de la scène mathématique, une
« stupéfiante \emph{dégradation} » de l'esprit du milieu, dont il se demande
si son départ n'a pas été le catalyseur, et qu'il dit n'avoir pas su percevoir
(« à la seule exception peut-être de Claude Chevalley ») ; il en isole un
aspect : le mépris de « ce qui ne s'apparente pas au pur travail du marteau sur
l'enclume » — l'inspiration, le rêve, la vision, et l'idée même tant qu'elle
n'est pas un énoncé publié, « répertoriable ». Il reprend à Siegel le mot
d'« aplatissement » de la pensée mathématique, privée de son « versant
d'ombre », du « versant féminin », et craint une « époque de dessèchement ». La
page commence et finit au milieu d'une phrase ; les renvois aux notes
(*) à (****) visent des notes absentes du dossier.

\end{document}
